\documentclass[11pt,twoside,fleqn]{article}

\usepackage[letterpaper,margin=0.68in]{geometry}
\usepackage[T1]{fontenc}
\usepackage{microtype}
\usepackage{amsmath,amssymb,amsthm,bm}
\usepackage{graphicx}
\usepackage{authblk}
\usepackage[numbers,sort&compress]{natbib}
\usepackage[colorlinks=true,linkcolor=blue,citecolor=blue,urlcolor=blue]{hyperref}

\setlength{\columnsep}{0.25in}
\setlength{\parskip}{0pt}
\setlength{\parindent}{1em}
\setlength{\mathindent}{0pt}

\renewcommand{\d}{\mathrm{d}}

\title{\textbf{Conditional Score-Based Modeling of Effective Langevin Dynamics}}
\author[1]{Ludovico T. Giorgini\thanks{ludogio@mit.edu; \url{https://ludogiorgi.github.io/}}}
\affil[1]{Department of Mathematics, Massachusetts Institute of Technology, Cambridge, MA 02139, USA}
\date{}

\begin{document}

\makeatletter
\twocolumn[
\begin{@twocolumnfalse}
\maketitle

\begin{abstract}
Stochastic reduced-order models are widely used to represent the effective
dynamics of complex systems, but estimating their drift and diffusion
coefficients from data remains challenging. Standard approaches often rely on
short-time trajectory increments, state-space partitioning, or repeated
simulation of candidate models, which become unreliable or computationally
expensive for high-dimensional systems, coarse temporal sampling, or unevenly
sampled data. We introduce a data-driven calibration method based on a novel
relationship between the coefficients of a stochastic reduced model and the
conditional score of the finite-time transition density, defined as the gradient
of the logarithm of the transition density with respect to the initial state.
The resulting identity expresses derivatives of lagged correlation functions as
stationary expectations over observed lagged pairs involving this conditional
score and the unknown model coefficients, or equivalently as finite-lag GFDT
impulse responses. This formulation allows the drift and diffusion structure to
be constrained
directly from finite-lag statistics, without differentiating trajectories,
partitioning state space, or repeatedly integrating candidate reduced models during
calibration, yielding a least-squares fitting problem over stationary lagged
pairs. We validate the approach on analytically tractable and
data-driven nonequilibrium diffusions, demonstrating that the inferred models
preserve the invariant statistics while accurately reproducing finite-lag
dynamical correlations. The framework provides a scalable route for learning
stochastic reduced-order models from data that reproduce prescribed statistical
and dynamical properties.
\end{abstract}

\vspace{1em}
\end{@twocolumnfalse}
]
\makeatother

\section{Introduction}

A central challenge in the study of high-dimensional multiscale dynamical systems is to construct reduced descriptions that retain the statistical and dynamical influence of unresolved variables. In many applications, the resolved variables represent only a small subset of the full state, while the eliminated degrees of freedom continue to affect their evolution through fast fluctuations, memory effects, and effective transport mechanisms. Projection and coarse-graining formalisms provide a precise framework for this idea: eliminating variables from a high-dimensional dynamics generally produces reduced equations with nonlocal memory terms and fluctuating forcing, even when the original dynamics are deterministic or Markovian in the complete state space \cite{Zwanzig1961,Mori1965,ChorinHaldKupferman2000,GivonKupfermanStuart2004,WoutersLucarini2012,WoutersLucarini2013,pavliotis2008multiscale,MTV1,MTV2}. In practice, one often replaces this non-Markovian reduced dynamics by an explicit Markovian stochastic surrogate, typically an effective stochastic differential equation whose drift and diffusion terms summarize the influence of unresolved variables on the observed state. Such stochastic parameterizations provide tractable and interpretable models for the statistics and dynamics of resolved observables, and have played a central role in climate dynamics, turbulence, and empirical model reduction \cite{Hasselmann1976,Penland1989,PenlandSardeshmukh1995,KONDRASHOV201533,ChekrounSimonnetGhil2011,LucariniChekroun2023}.

A stochastic reduced-order model should preserve the statistical and dynamical features of the observed variables that are relevant at the resolved scale. A necessary requirement is the reproduction of the invariant measure, since the stationary distribution encodes the long-time statistics of the resolved observables. This requirement alone, however, is not sufficient. Distinct stochastic dynamics can share the same invariant density while exhibiting different temporal correlations, probability currents, response properties, and effective dynamical couplings. A successful reduced model must therefore reproduce not only the steady-state distribution, but also the dynamical mechanisms by which probability is transported through state space. This distinction is especially important in climate modeling, where one often seeks to predict how coarse observables respond to perturbations, rather than only to reproduce their unforced climatology. Response theory makes clear that forced responses are controlled by the interplay between invariant statistics and dynamical transport, not by the invariant density alone \cite{Ruelle1998,Ruelle2009}. Recent analyses of data-driven climate emulators and reduced models have reinforced this point, showing that accurate stationary statistics do not by themselves guarantee accurate response behavior or physically meaningful coarse dynamics \cite{falasca2025probing,falasca2024data,falasca2026causalrom,lucarini2026koopmanism}. The construction of stochastic reduced models should therefore be formulated as the joint problem of matching invariant statistics and finite-lag dynamical observables.

A first route is calibration by repeated simulation: one specifies a parametric or functional model ansatz, integrates the resulting dynamics for each trial parameter value, estimates from the simulated trajectory the statistical or dynamical observables of interest, compares them with the corresponding targets, and updates the model accordingly. This strategy is closely related to simulation-based inference approaches, in which model parameters are inferred by matching statistics or discrepancies computed from simulated and observed data \cite{DuffieSingleton1993,GourierouxMonfortRenault1993,GutmannCorander2016}. Although very general, such approaches become computationally inefficient when each model evaluation requires a long integration of a high-dimensional stochastic surrogate. A second route is to infer the local-in-time dynamical law directly from trajectory data. This includes Kramers--Moyal and Fokker--Planck reconstruction methods, which estimate drift and diffusion coefficients from short-time conditional increment statistics, as well as related stochastic-modeling approaches that recover effective tendencies, stability coefficients, or noise amplitudes from finite differences of observed trajectories \cite{Friedrich1997,Siegert1998,RagwitzKantz2001,PenlandSardeshmukh1995,giorgini2022non,keyes2023stochastic}. Such methods have provided important tools for empirical stochastic modeling, but they are intrinsically tied to the availability of sufficiently resolved time series: when observations are noisy, unevenly sampled, or separated by time intervals much longer than the intrinsic dynamical time scale, finite-difference estimates of local tendencies and short-time conditional moments become unreliable. A third route is to approximate coarse generators, transfer operators, or lagged correlation structures through state-space partitions, clustering, or related operator-theoretic constructions \cite{GivonKupfermanStuart2004,BotvinickGreenhouse2023,BotvinickGreenhouse2025,SantosGutierrezLucariniChekrounGhil2021,falasca2024data,falasca2025FDT,giorgini2025learning,giorgini2024reduced,souza2024representing_a,souza2024representing_b,souza2024modified}. These approaches are powerful and conceptually close to the present work, but direct reconstruction of local drift and diffusion from partitions suffers from the curse of dimensionality. These limitations motivate an inverse formulation in which an effective stochastic generator is constrained directly from stationary samples and lagged pairs, without repeated forward integration during calibration, without estimating instantaneous rates of change from trajectories, and without reconstructing local coefficients by state-space clustering.

The alternative pursued here is to use finite-lag conditional statistics as the primary dynamical object. This viewpoint is motivated by time-delay embedding theory, which shows that, for deterministic systems and under suitable generic
conditions, lagged observations can encode the dynamics on the underlying
attractor \cite{Takens1981,BotvinickGreenhouse2025}. In the stochastic
multiscale setting considered here, however, we do not rely on a deterministic
embedding theorem. Rather, the corresponding object is the finite-time
transition law of the resolved variables, which describes how the present
resolved state statistically constrains future resolved states, even when the
full high-dimensional state is unobserved. The relevant object is therefore not
an instantaneous drift or local tendency estimated from consecutive data points,
but the conditional law of future resolved states given present resolved states.
We represent this conditional law through its score, namely the gradient of the
logarithm of the finite-lag transition density with respect to the present
resolved state. The conditional score provides a local sensitivity
representation of the finite-lag statistical dependence between present and
future states, and thus contains the information needed to constrain the reduced
generator. Recent advances in score matching and score-based generative modeling
make it possible to estimate such score functions directly from high-dimensional
samples without explicit density reconstruction
\cite{Hyvarinen2005,Vincent2011,SongSohlDicksteinKingma2021,giorgini2025kgmm}.
The purpose of this paper is to derive a direct relation between conditional
scores and the mobility structure of an effective Langevin model, and to use
this relation to infer stochastic reduced dynamics from stationary samples and
lagged pairs, avoiding the limitations of simulation-based calibration,
state-space partitioning, and trajectory differentiation.

The use of score-function estimators developed in generative score modeling has already proved highly effective in nonequilibrium statistical physics and data-driven modeling of nonlinear dynamical systems. One important application is generalized fluctuation--dissipation theory, which relates the response of observables to weak perturbations to correlation functions evaluated in the unperturbed system \cite{Kubo1966FDT,marconi2008,Ruelle1998,Ruelle2009,majda2009,majdaStructuralStability,MajdaBook,GRITSUN,cooper2011climate,LucariniChekroun2023}. In nonlinear nonequilibrium systems, these correlation functions involve derivatives of the invariant density, a quantity that is difficult to estimate directly in high dimension. Learned score functions provide precisely this information, and have enabled accurate data-driven response predictions for nonlinear and high-dimensional dynamical systems \cite{giorgini_response_theory,giorgini2025probdist,giorgini2025calibration}. The same principle has also been used to construct score-based Langevin reduced models from data \cite{giorgini2026scorelangevin, del2025integrating}. In that framework, the learned stationary score fixes the target steady-state density, while a constant mobility matrix, determined from the short-time derivative of the coordinate autocorrelation, closes the dynamics by replacing the state-dependent transport structure with a mean-field approximation. This construction preserves the steady-state distribution by design and reproduces the leading short-time coordinate correlations, and has proved effective in high-dimensional examples, including the Kuramoto--Sivashinsky equation and cyclostationary sea-surface-temperature data from PlaSim \cite{giorgini2025cyclostationary}. Its limitation, however, is intrinsic to the closure: because the mobility is constant and only the stationary score is used, the model can enforce only the short-time correlation constraints used to determine that constant matrix, and cannot systematically impose more general finite-lag dynamical observables. The present work removes this restriction by introducing the conditional score as the object that links finite-lag transition statistics to the state-dependent mobility of the reduced Langevin generator. This extension turns score-based stochastic modeling from a construction that preserves the invariant density and a restricted set of short-time correlations into a framework in which prescribed dynamical constraints can be imposed directly from data.

The remainder of the paper is organized as follows.
Section~\ref{sec:method} introduces the score-based Langevin representation,
derives the central finite-lag conditional-score identity, and formulates the
inverse problem for the mobility.
Section~\ref{sec:results} presents analytically tractable and data-driven
examples, and Section~\ref{sec:conclusions} summarizes the main implications.
Technical derivations and implementation details are collected in the
Appendices.

\section{Method}
\label{sec:method}

\subsection{Problem setting and modeling objective}
\label{subsec:problem_setting_modeling_objective}

We consider a resolved process $\bm{x}(t)\in\mathbb{R}^D$ obtained from observations. The goal is to construct, directly from data, an effective stochastic model for these resolved variables. The construction is based on four assumptions.

First, the resolved process is assumed to admit a statistically stationary distribution $p_{\mathrm{ss}}(\bm{x})$. Cyclostationary data can be incorporated by augmenting the state with the phase of the external forcing, while slow nonstationary trends should be removed before applying the method. Second, the process is assumed to be ergodic and sufficiently mixing. Ergodicity is needed to replace ensemble expectations by time averages along a single sufficiently long realization, while mixing ensures that correlations between widely separated samples decay, which stabilizes the empirical lagged statistics entering the inverse problem. Third, the resolved variables are assumed to evolve on slower time scales than the unresolved degrees of freedom, which appear only through their effective stochastic influence on the resolved dynamics. We assume, in the main presentation, that the slow variables are fully observed; if the observations are partial, the same framework can be applied after replacing the state by a delay-embedded state. In that setting, the delay coordinates carry finite-window information about the unobserved components and their memory effects, so that the effective reduced dynamics can again be modeled on an augmented observed state. Finally, the data must have adequate temporal resolution for the dynamical quantities being matched. In particular, we do not require the sampling to be fine enough to estimate instantaneous time derivatives, but it must resolve the finite-lag statistics that will be used as modeling constraints.

We model the resolved dynamics by an It\^o diffusion of the form
\begin{equation}
    \d \bm{x}(t)
    =
    \bm{F}(\bm{x}(t))\,\d t
    +
    \sqrt{2}\,\bm{\Sigma}(\bm{x}(t))\,\d \bm{W}_t,
    \label{eq:method_sde_full}
\end{equation}
where $\bm{F}:\mathbb{R}^D\to\mathbb{R}^D$ is the drift,
$\bm{\Sigma}:\mathbb{R}^D\to\mathbb{R}^{D\times D}$ is the noise-amplitude matrix, and $\bm{W}_t$ is a standard $D$-dimensional Wiener process. We denote the diffusion tensor by
\begin{equation}
    \bm{D}(\bm{x})
    :=
    \bm{\Sigma}(\bm{x})\bm{\Sigma}(\bm{x})^T .
    \label{eq:diffusion_tensor_definition}
\end{equation}
The central modeling problem is to infer $\bm{F}$ and $\bm{D}$ from data in such a way that the resulting diffusion reproduces both the stationary distribution of the observations and selected dynamical features of the resolved process.

To enforce the stationary distribution, we use the score of the observed steady state,
\begin{equation}
    \bm{s}(\bm{x})
    :=
    \nabla \log p_{\mathrm{ss}}(\bm{x}).
    \label{eq:score_definition_method}
\end{equation}
The stationary score can be estimated directly from samples by denoising score matching; see Appendices~\ref{app:stationary_score_dsm} and~\ref{app:conditional_score_dsm} for the stationary and conditional constructions used throughout the paper.
Following the score-based Langevin representation derived in \cite{giorgini2026scorelangevin}, a stationary diffusion with invariant density $p_{\mathrm{ss}}$ can be written in the score-based form below. A detailed derivation and the required regularity and boundary assumptions are given in Appendix~\ref{app:mobility_representation}.
\begin{equation}
    \d \bm{x}(t)
    =
    \left[
        \bm{M}(\bm{x})\bm{s}(\bm{x})
        +
        \nabla\cdot\bm{M}(\bm{x})
    \right]\d t
    +
    \sqrt{2}\,\bm{\Sigma}(\bm{x})\,\d \bm{W}_t,
    \label{eq:compact_sde_full}
\end{equation}
where
\begin{equation}
\begin{aligned}
    \bm{M}(\bm{x})
    &=
    \bm{D}(\bm{x})
    +
    \bm{R}(\bm{x}),
    \\
    \bm{R}(\bm{x})^T
    &=
    -\bm{R}(\bm{x}),
    \\
    \bm{D}(\bm{x})
    &=
    \frac{\bm{M}(\bm{x})+\bm{M}(\bm{x})^T}{2}
    \succeq 0 .
\end{aligned}
    \label{eq:mobility_decomposition_method}
\end{equation}
Here $(\nabla\cdot\bm{M})_i=\sum_{j=1}^D\partial_{x_j}M_{ij}$. The symmetric part of $\bm{M}$ determines the diffusion tensor, while the antisymmetric part represents stationary probability currents and therefore allows the reduced model to describe nonequilibrium dynamics. The advantage of \eqref{eq:compact_sde_full} is that the steady-state density is imposed by construction: once the score $\bm{s}$ and the mobility $\bm{M}$ are specified, the model has $p_{\mathrm{ss}}$ as invariant density.

Matching the stationary density alone, however, is not sufficient. Many stochastic processes can share the same invariant distribution while having very different temporal organization, response properties, and effective transport. We therefore constrain the mobility using lagged statistics of a prescribed family of observables. Let
\begin{equation}
    \phi_m:\mathbb{R}^D\to\mathbb{R}^{d_m},
    \qquad
    m=1,\dots,M,
    \label{eq:observable_family_method}
\end{equation}
be smooth observables chosen to represent the dynamical features of interest. For each pair $(m,n)$ and each lag $t$ in a selected set of resolved time scales, we define the observed lagged correlation matrix
\begin{equation}
    \bm{C}_{m,n,\mathrm{obs}}(t)
    :=
    \left\langle
        \phi_m(\bm{x}(t))\,
        \phi_n(\bm{x}(0))^T
    \right\rangle_{\mathrm{obs}} .
    \label{eq:Cm_definition}
\end{equation}
The modeling objective is then to construct a diffusion of the form \eqref{eq:compact_sde_full} such that
\begin{equation}
\begin{aligned}
    p_{\mathrm{model}}(\bm{x})
    &\approx
    p_{\mathrm{ss}}(\bm{x}),
    \\
    \bm{C}_{m,n,\mathrm{model}}(t)
    &\approx
    \bm{C}_{m,n,\mathrm{obs}}(t).
\end{aligned}
    \label{eq:modeling_objective}
\end{equation}
The lagged-correlation matching is imposed for the selected observables and lags.

This choice of constraints is deliberate. The invariant density fixes the statistics of the reduced process, while the lagged correlations probe the action of the dynamics on observables and therefore encode information about time scales, transport, memory effects, and linear response. At the same time, these quantities are statistically robust: they can be estimated as time averages from observed trajectories, including sparsely or unevenly sampled data by using all available pairs within prescribed lag windows. The remainder of this section shows how the mobility $\bm{M}$ can be inferred from the score and from the empirical lagged correlations, without requiring repeated long integrations of candidate reduced models.

\subsection{Inferring the mobility from correlation constraints}
\label{subsec:mobility_from_correlation_constraints}

The preceding subsection reduces the modeling problem to the identification of the mobility field $\bm{M}(\bm{x})$. The steady-state score fixes the invariant density. The mobility then determines the probability currents, time scales, transport properties, and lagged correlations of the reduced process. We now show how $\bm{M}$ can be inferred directly from the empirical correlation constraints introduced above. The conditional scores entering this construction are estimated from lagged pairs by the joint-score denoising score-matching procedure summarized in Appendix~\ref{app:conditional_score_dsm}.
The key object is the finite-lag transition density of the observed resolved process,
\begin{equation}
    p_t(\bm{x}_t\mid \bm{x}_0),
\end{equation}
and its score with respect to the initial condition,
\begin{equation}
    \bm{s}_{t|0}(\bm{x}_t\mid \bm{x}_0)
    :=
    \nabla_{\bm{x}_0}
    \log p_t(\bm{x}_t\mid \bm{x}_0).
    \label{eq:conditional_score_method_section}
\end{equation}
This conditional score measures how the likelihood of observing $\bm{x}_t$ at lag $t$ changes under an infinitesimal perturbation of the initial state $\bm{x}_0$. It is therefore the finite-time analogue of the stationary score, but applied to transition probabilities rather than to the invariant density.

The central identity of the method is that, for the score-based diffusion
\eqref{eq:compact_sde_full}, the derivative of the lagged correlation matrix can
be written in the equivalent forms
\begin{equation}
\begin{aligned}
    \dot{\bm{C}}_{m,n}(t)
    &=
    -
    \left\langle
        \nabla (K_t\phi_m)(\bm{x}(0))\,
        \bm{M}(\bm{x}(0))
    \right. \\
    &\qquad\left.
        \nabla\phi_n(\bm{x}(0))^T
    \right\rangle
    \\
    &=
    -
    \left\langle
        \phi_m(\bm{x}(t))\,
        \bm{s}_{t|0}(\bm{x}(t)\mid \bm{x}(0))^T
    \right. \\
    &\qquad\left.
        \bm{M}(\bm{x}(0))\,
        \nabla\phi_n(\bm{x}(0))^T
    \right\rangle .
\end{aligned}
    \label{eq:central_conditional_score_identity}
\end{equation}
Here \(K_t\) denotes the Koopman semigroup associated with the reduced
Markov process, \(\nabla(K_t\phi_m)\) is the Jacobian of the lag-\(t\)
conditional expectation of the observable \(\phi_m\), and \(\nabla\phi_n\)
denotes the Jacobian of the vector-valued observable \(\phi_n\). The
expectation is over stationary lagged pairs \((\bm{x}(0),\bm{x}(t))\).
The second equality follows from the transition-density representation of
\(K_t\phi_m\) and from the definition of the conditional score
\(\bm{s}_{t|0}=\nabla_{\bm{x}(0)}\log p_t(\bm{x}(t)\mid\bm{x}(0))\). A
detailed derivation of \eqref{eq:central_conditional_score_identity} is given in
Appendix~\ref{app:conditional_score_correlation_identity}.

The first line of \eqref{eq:central_conditional_score_identity} makes explicit
the Koopman-theoretic structure of the inverse problem. In Koopman spectral
approaches, one typically represents the action of \(K_t\) through eigenvalues
and eigenfunctions, so that, formally, if
\(\phi_m=\sum_{\ell}\bm{a}_{m,\ell}\psi_\ell\) and
\(K_t\psi_\ell=e^{\lambda_\ell t}\psi_\ell\), then
\[
    \nabla(K_t\phi_m)(\bm{x})
    =
    \sum_{\ell}
        e^{\lambda_\ell t}\,
        \bm{a}_{m,\ell}\,
        \nabla\psi_\ell(\bm{x}) .
\]
This viewpoint underlies Koopman mode decompositions, dynamic mode
decomposition, extended dynamic mode decomposition, and recent formulations
linking Koopman spectral decompositions with response theory
\cite{Koopman1931,Mezic2005,BudisicMohrMezic2012,WilliamsKevrekidisRowley2015,KlusNuskeKoltaiWuKevrekidisSchutteNoe2018,SantosGutierrezLucariniChekrounGhil2021,ZagliColbrookLucariniMezicMoroney2026,LucariniSantosGutierrezMoroneyZagli2026}.
In that setting, response and correlation operators can be decomposed into
contributions associated with Koopman modes, providing a highly interpretable
description of the dynamical time scales.

For high-dimensional stochastic or coarse-grained systems, however, using this
representation as a practical inverse tool requires estimating sufficiently many
Koopman eigenfunctions and, in the present identity, their spatial gradients.
This is a severe numerical bottleneck: the eigenfunction approximation depends
on the choice of dictionary, kernel, or finite-rank subspace; the required number
of modes can grow rapidly with dimension; and differentiating the learned
eigenfunctions can amplify approximation error. The second line of
\eqref{eq:central_conditional_score_identity} provides an alternative route. It
retains the same Koopman object, \(\nabla(K_t\phi_m)\), but evaluates it through
the conditional transition score instead of through an explicit spectral
resolution of the Koopman operator. Thus, neural score estimation on lagged
pairs supplies the differential information needed by the Koopman-gradient
identity without constructing Koopman eigenfunctions. This is the key practical
advantage of the present formulation: it turns a Koopman response identity into
a data-driven correlation constraint that remains linear in the unknown mobility
\(\bm{M}\), avoids instantaneous velocity estimation, and can be implemented in
high-dimensional systems using modern conditional score estimators.

In practice, we choose a set of resolved lags $\mathcal{T}$ and replace all expectations in \eqref{eq:central_conditional_score_identity} by empirical averages over observed pairs. This defines the empirical correlation operator
\begin{equation}
\begin{aligned}
    \widehat{\mathcal{A}}_{m,n,t}[\bm{M}]
    &:=
    -
    \left\langle
        \phi_m(\bm{x}_t)\,
        \widehat{\bm{s}}_{t|0}(\bm{x}_t\mid \bm{x}_0)^T
    \right. \\
    &\qquad\left.
        \bm{M}(\bm{x}_0)\,
        \nabla\phi_n(\bm{x}_0)^T
    \right\rangle_{\mathrm{obs}},
    \qquad
    t\in\mathcal{T}.
\end{aligned}
    \label{eq:empirical_mobility_operator}
\end{equation}
The mobility is then determined by enforcing
\begin{equation}
\begin{aligned}
    \widehat{\mathcal{A}}_{m,n,t}[\bm{M}]
    &\approx
    \dot{\bm{C}}_{m,n,\mathrm{obs}}(t),
    \\
    &\qquad
    t\in\mathcal{T},
    \quad
    m,n=1,\dots,M_{\mathrm{obs}},
\end{aligned}
    \label{eq:mobility_correlation_constraint}
\end{equation}
together with the structural constraint
\begin{equation}
    \frac{\bm{M}(\bm{x})+\bm{M}(\bm{x})^T}{2}\succeq 0 .
    \label{eq:mobility_positive_constraint}
\end{equation}
For a finite-dimensional parametrization $\bm{M}_\theta$, this leads to the least-squares problem
\begin{equation}
\begin{aligned}
    \mathcal{J}(\theta)
    &=
    \sum_{t\in\mathcal{T}}
    \sum_{m,n}
    w_{m,n,t}
    \left\|
        \widehat{\mathcal{A}}_{m,n,t}[\bm{M}_\theta]
        -
        \dot{\bm{C}}_{m,n,\mathrm{obs}}(t)
    \right\|_F^2
    \\
    &\quad+
    \lambda\,\mathcal{R}(\theta),
\end{aligned}
    \label{eq:mobility_objective_general}
\end{equation}
where $\mathcal{R}$ denotes regularization. The weights $w_{m,n,t}$ may account for sampling uncertainty or for the relative importance of different observables and lags.

The reason for matching correlation derivatives, rather than only the correlations themselves, is that the derivative constraint identifies the generator of the reduced process on the chosen observable class. Since the invariant density has already been imposed by construction, the model and the observations have the same equal-time correlations:
\begin{equation}
    \bm{C}_{m,n,\mathrm{model}}(0)
    =
    \bm{C}_{m,n,\mathrm{obs}}(0).
    \label{eq:correlation_initial_match}
\end{equation}
Therefore, for any lag $t$ in the resolved interval,
\begin{equation}
\begin{aligned}
    \bm{C}_{m,n,\mathrm{model}}(t)
    -
    \bm{C}_{m,n,\mathrm{obs}}(t)
    &=
    \int_0^t
    \Big[
        \dot{\bm{C}}_{m,n,\mathrm{model}}(\tau)
        \\
    &\qquad-
        \dot{\bm{C}}_{m,n,\mathrm{obs}}(\tau)
    \Big]\d \tau .
\end{aligned}
    \label{eq:correlation_error_integral}
\end{equation}
Thus, if the mobility inferred from \eqref{eq:mobility_correlation_constraint} makes the model correlation derivatives agree with the observed derivatives over the selected lag range, then the corresponding correlation functions agree over that range. More precisely, if
\begin{equation}
\begin{aligned}
    \left\|
        \dot{\bm{C}}_{m,n,\mathrm{model}}(\tau)
        -
        \dot{\bm{C}}_{m,n,\mathrm{obs}}(\tau)
    \right\|_F
    &\leq
    \varepsilon_{m,n}(\tau),
    \\
    &\qquad
    0\leq \tau\leq T,
\end{aligned}
    \label{eq:correlation_derivative_error_bound}
\end{equation}
then
\begin{equation}
    \left\|
        \bm{C}_{m,n,\mathrm{model}}(t)
        -
        \bm{C}_{m,n,\mathrm{obs}}(t)
    \right\|_F
    \leq
    \int_0^t
        \varepsilon_{m,n}(\tau)\,\d \tau .
    \label{eq:correlation_error_bound}
\end{equation}
The derivative constraint is therefore a weak, data-driven way of enforcing dynamical agreement on the observable library.

We now clarify what it means for the observable library to be sufficiently rich. A strong sufficient
condition is that the selected observables and lags identify the admissible
mobility perturbations. Let $\mathcal{H}$ denote the admissible class of
mobility perturbations, for example the tangent space associated with the
chosen parametrization of $\bm{M}$. Define the nullspace
\begin{equation}
\begin{aligned}
    &\mathcal{N}_{\mathcal{T},\phi}(\mathcal{H})
    :=
    \Bigl\{
        \bm{H}\in\mathcal{H}: \\
        &\quad\left\langle
            \phi_m(\bm{x}_t)\,
            \bm{s}_{t|0}(\bm{x}_t\mid\bm{x}_0)^T
            \bm{H}(\bm{x}_0)\,
            \nabla\phi_n(\bm{x}_0)^T
        \right\rangle_{\mathrm{obs}}
        =0, \\
        &\quad\text{for all } t\in\mathcal{T}, \text{ and } m,n
    \Bigr\}.
\end{aligned}
    \label{eq:mobility_operator_nullspace}
\end{equation}
If
\begin{equation}
\begin{aligned}
    \mathcal{N}_{\mathcal{T},\phi}
    \bigl(\mathcal{H}\bigr)
    &=
    \{\bm{0}\},
\end{aligned}
    \label{eq:restricted_injectivity_condition_method}
\end{equation}
then the selected constraints are separating on $\mathcal{H}$: no nonzero
admissible perturbation of the mobility leaves all selected
correlation-derivative constraints unchanged. This restricted-injectivity
condition is therefore sufficient for identifying the mobility on the support
of the invariant measure, up to the resolution of the chosen parametrization.

For the reduced-modeling objective pursued here, however, this condition is
stronger than necessary. The practical requirement is that the observable
library be rich enough that a mobility satisfying the empirical GFDT constraints
\begin{equation}
\begin{aligned}
    \dot C_{mn}^{\mathrm{obs}}(t)
    &\approx
    -
    \left\langle
        \phi_m(\bm{x}_t)\,
        \bm{s}_{t|0}^{\mathrm{obs}}(\bm{x}_t\mid\bm{x}_0)^T
    \right. \\
    &\qquad\left.
        \bm{M}(\bm{x}_0)
        \nabla\phi_n(\bm{x}_0)^T
    \right\rangle_{\mathrm{obs}} .
\end{aligned}
    \label{eq:weak_observable_adequacy}
\end{equation}
also yields a model whose corresponding correlation derivatives are close to
the observed ones,
\begin{equation}
    \dot C_{mn}^{\mathrm{model}}(t)
    \approx
    \dot C_{mn}^{\mathrm{obs}}(t),
    \qquad
    t\in\mathcal{T}.
    \label{eq:model_correlation_derivative_adequacy}
\end{equation}
Thus, strict injectivity is an identifiability condition. The weaker condition
needed in practice is observable-level adequacy. In this weaker setting,
different mobilities may be indistinguishable by the selected observables and
lags. That nonuniqueness is harmless if all such mobilities produce the same
resolved correlation dynamics. With finite data, the relevant criterion is not
exact nullspace triviality. It is stability and predictive closure: the
selected constraints should make the inverse problem well conditioned for the
correlation statistics one aims to reproduce.

This is the sense in which the method avoids both local velocity reconstruction
and state-space clustering. The dynamical information is extracted from robust
finite-lag statistics and from conditional scores of observed pairs. The
mobility is then inferred within the chosen admissible class by matching the
weak correlation dynamics of the data on a selected family of observables.

\subsection{GFDT interpretation of the constraint}
\label{subsec:gfdt_interpretation_correlation_constraint}

The correlation identity derived above can also be viewed as a generalized
fluctuation--dissipation relation. This perspective is useful because it
identifies the correlation-derivative constraint as an impulse-response kernel
of the stationary reduced dynamics. We use the same regularity and boundary
assumptions as in Appendix~\ref{app:conditional_score_correlation_identity}, so
that integrations by parts against the invariant density are justified.

Consider the unperturbed score-based diffusion \eqref{eq:compact_sde_full},
with invariant density \(p_{\mathrm{ss}}\), and perturb its drift by a small
time-dependent vector field. For a perturbing drift
\(\bm{b}:\mathbb{R}^D\to\mathbb{R}^D\), the perturbed dynamics are
\begin{align}
    \d\bm{x}_\varepsilon(t)
    &=
    \left[
        \bm{M}(\bm{x}_\varepsilon)\bm{s}(\bm{x}_\varepsilon)
        +
        \nabla\cdot\bm{M}(\bm{x}_\varepsilon)
    \right]\d t
    \nonumber
    \\
    &\quad+
    \varepsilon h(t)\bm{b}(\bm{x}_\varepsilon)\,\d t
    \nonumber
    \\
    &\quad+
    \sqrt{2}\,\bm{\Sigma}(\bm{x}_\varepsilon)\,\d\bm{W}_t .
    \label{eq:gfdt_perturbed_sde}
\end{align}
The generalized fluctuation--dissipation theorem gives, to first order in
\(\varepsilon\), the response of a smooth observable
\(\bm{O}:\mathbb{R}^D\to\mathbb{R}^{d_O}\) as
\cite{Agarwal1972,FalcioniIsolaVulpiani1990,MarconiPuglisiRondoniVulpiani2008,Lucarini2008}
\begin{equation}
    \delta\left\langle \bm{O}(\bm{x}(t))\right\rangle
    =
    \varepsilon
    \int_0^t
        h(\tau)\,
        \bm{R}_{\bm{O},\bm{b}}(t-\tau)\,
    \d\tau
    +
    o(\varepsilon),
    \label{eq:gfdt_convolution_response}
\end{equation}
where the response kernel is the stationary correlation
\begin{align}
    \bm{R}_{\bm{O},\bm{b}}(t)
    &=
    \left\langle
        \bm{O}(\bm{x}(t))\,
        B_{\bm{b}}(\bm{x}(0))
    \right\rangle ,
    \\
    B_{\bm{b}}(\bm{x})
    &=-
    \frac{1}{p_{\mathrm{ss}}(\bm{x})}
    \nabla\cdot
    \left[
        p_{\mathrm{ss}}(\bm{x})\bm{b}(\bm{x})
    \right].
    \label{eq:gfdt_conjugate_variable_vector}
\end{align}
For an impulsive perturbation, \(h(t)=\delta(t)\), the response is therefore
directly the correlation
\(\bm{R}_{\bm{O},\bm{b}}(t)\).

We now apply this formula to perturbations indexed by the components of
\(\phi_n\). For a matrix field \(\bm{A}(\bm{x})\), define the componentwise
drift perturbations
\begin{equation}
    \bm{b}_{\bm{A},n,j}(\bm{x})
    :=
    -
    \bm{A}(\bm{x})\nabla\phi_{n,j}(\bm{x}),
    \qquad
    j=1,\dots,d_n .
    \label{eq:b_A_n_j_definition}
\end{equation}
The corresponding conjugate variables are the scalar observables
\begin{equation}
    B_{\bm{b}_{\bm{A},n,j}}(\bm{x})
    =
    \frac{1}{p_{\mathrm{ss}}(\bm{x})}
    \nabla\cdot
    \left[
        p_{\mathrm{ss}}(\bm{x})\,
        \bm{A}(\bm{x})\nabla\phi_{n,j}(\bm{x})
    \right].
    \label{eq:B_b_A_n_j_definition}
\end{equation}

For the physical mobility \(\bm{A}=\bm{M}\), the \(j\)-th response column is
the response to the drift impulse \(\bm{b}_{\bm{M},n,j}\):
\begin{equation}
    \left[\bm{\mathcal{R}}_{\phi_m,\bm{M},n}(t)\right]_{:,j}
    =
    \left\langle
        \phi_m(\bm{x}(t))\,
        B_{\bm{b}_{\bm{M},n,j}}(\bm{x}(0))
    \right\rangle .
    \label{eq:gfdt_response_M_grad_phi_n}
\end{equation}
This response matrix is exactly the derivative of the lagged
correlation \(\bm{C}_{m,n}(t)\). Using the Koopman semigroup and then
integrating by parts against \(p_{\mathrm{ss}}\), for each component \(j\),
\begin{align}
    &\left\langle
        \phi_m(\bm{x}(t))\,
        B_{\bm{b}_{\bm{M},n,j}}(\bm{x}(0))
    \right\rangle
    \nonumber\\
    &\qquad=
    \left\langle
        (K_t\phi_m)(\bm{x}_0)\,
        B_{\bm{b}_{\bm{M},n,j}}(\bm{x}_0)
    \right\rangle_{p_{\mathrm{ss}}}
    \nonumber\\
    &\qquad=
    -
    \left\langle
        \nabla(K_t\phi_m)(\bm{x}_0)\,
        \bm{M}(\bm{x}_0)\,
        \nabla\phi_{n,j}(\bm{x}_0)
    \right\rangle_{p_{\mathrm{ss}}}.
    \label{eq:gfdt_response_integration_by_parts}
\end{align}
The final expression is precisely the \(j\)-th column of the gradient form of
\eqref{eq:central_conditional_score_identity}. Hence
\begin{equation}
    \dot{\bm{C}}_{m,n}(t)
    =
    \bm{\mathcal{R}}_{\phi_m,\bm{M},n}(t).
    \label{eq:Cdot_as_gfdt_response}
\end{equation}
Equivalently, the correlation-derivative constraint is the collection of GFDT
impulse responses generated by the componentwise drift fields
\(\bm{b}_{\bm{M},n,j}\).


In the reversible equilibrium limit, the general identity reduces to the
classical fluctuation--dissipation theorem
\cite{CallenWelton1951,Kubo1957,Kubo1966}. Define the equilibrium correlation
\begin{equation}
    \bm{C}^{\mathrm{eq}}_{m,n}(t)
    :=
    \left\langle
        \phi_m(\bm{x}(t))\,
        \phi_n(\bm{x}(0))^T
    \right\rangle_{\mathrm{eq}} .
    \label{eq:equilibrium_correlation_definition}
\end{equation}
For a field coupled as \(U_{\bm{h}}=U-\bm{h}^T\phi_n\), the equilibrium impulse
response satisfies
\begin{equation}
    \bm{\mathcal{R}}^{\mathrm{eq}}_{m,n}(t)
    =
    -\beta\,\frac{\d}{\d t}
    \bm{C}^{\mathrm{eq}}_{m,n}(t),
    \qquad t>0.
    \label{eq:classical_equilibrium_fdt}
\end{equation}
Here \(T\) is the absolute temperature, \(k_{\mathrm B}\) is Boltzmann's
constant, and
\begin{equation}
    \beta=(k_{\mathrm B}T)^{-1}
    \label{eq:inverse_temperature_definition}
\end{equation}
is the inverse thermal energy. The equilibrium limit corresponds to setting the
antisymmetric mobility component to zero, \(\bm{R}=0\), and taking the
diffusion tensor to be spatially uniform and isotropic,
\begin{equation}
    \bm{D}=\beta^{-1}\bm{I}.
    \label{eq:equilibrium_isotropic_diffusion}
\end{equation}
Thus the thermal noise strength is fixed by the temperature through the
Einstein relation: larger temperature corresponds to stronger diffusion.

In this case the stationary density can be written as a Boltzmann distribution.
Let \(U:\mathbb{R}^D\to\mathbb{R}\) be a confining potential and define the
partition function
\begin{equation}
    Z
    :=
    \int_{\mathbb{R}^D}
        e^{-\beta U(\bm{y})}
    \,\d\bm{y}
    <
    \infty .
    \label{eq:equilibrium_partition_function}
\end{equation}
Then
\begin{equation}
    p_{\mathrm{ss}}(\bm{x})
    =
    Z^{-1}e^{-\beta U(\bm{x})}.
    \label{eq:equilibrium_boltzmann_density}
\end{equation}
It follows that \(\bm{s}=\nabla\log p_{\mathrm{ss}}=-\beta\nabla U\), and the
dynamics become
\begin{equation}
\begin{aligned}
    \d\bm{x}(t)
    &=
    -\nabla U(\bm{x}(t))\,\d t
    +
    \sqrt{2\beta^{-1}}\,\d\bm{W}_t .
\end{aligned}
    \label{eq:equilibrium_overdamped_langevin}
\end{equation}
Moreover, \eqref{eq:Cdot_as_gfdt_response} gives the componentwise identity
\begin{equation}
    \bm{b}^{\mathrm{eq}}_{n,j}
    =
    -
    \beta^{-1}\nabla\phi_{n,j}.
    \label{eq:equilibrium_b_n_j}
\end{equation}
\begin{equation}
\begin{aligned}
    \left[\dot{\bm{C}}_{m,n}(t)\right]_{:,j}
    &=
    \left\langle
        \phi_m(\bm{x}(t))\,
        B_{\bm{b}^{\mathrm{eq}}_{n,j}}(\bm{x}(0))
    \right\rangle .
\end{aligned}
    \label{eq:equilibrium_cdot_b}
\end{equation}
Now perturb the equilibrium potential by coupling an external field
\(\bm{h}\) to \(\phi_n\),
\begin{equation}
    U_{\bm{h}}(\bm{x})
    =
    U(\bm{x})
    -
    \bm{h}^T\phi_n(\bm{x}) .
    \label{eq:equilibrium_field_coupling}
\end{equation}
The induced drift perturbation is
\begin{equation}
\begin{aligned}
    \delta\bm{F}_{\bm{h}}(\bm{x})
    &=
    \nabla\phi_n(\bm{x})^T\bm{h}
    \\
    &=
    \beta
    \left[
        \beta^{-1}\nabla\phi_n(\bm{x})^T
    \right]\bm{h}.
\end{aligned}
    \label{eq:equilibrium_field_drift_perturbation}
\end{equation}
The physical drift direction satisfies
\begin{equation}
    \nabla\phi_n^T\bm{h}
    =
    -\beta
    \sum_{j=1}^{d_n}
        h_j\bm{b}^{\mathrm{eq}}_{n,j}.
    \label{eq:equilibrium_physical_drift_b_decomposition}
\end{equation}
Therefore the response to this perturbation is
\begin{equation}
    \bm{\mathcal{R}}^{\mathrm{eq}}_{m,n}(t)
    =
    -\beta\,\dot{\bm{C}}_{m,n}(t).
    \label{eq:equilibrium_fdt_recovered}
\end{equation}
This is the standard equilibrium FDT for the coupling
\eqref{eq:equilibrium_field_coupling}.

\subsection{Mean mobility, mobility corrections, and the constant closure}
\label{subsec:mean_mobility_and_corrections}

The correlation constraints in Section~\ref{subsec:mobility_from_correlation_constraints} determine the mobility field only through its action on the chosen observable library. We introduce the decomposition
\begin{equation}
\begin{aligned}
    \bm{M}(\bm{x})
    &=
    \bm{\Phi}
    +
    \delta\bm{M}(\bm{x}),
    \\
    \bm{\Phi}
    &:=
    \left\langle
        \bm{M}
    \right\rangle_{\mathrm{obs}},
    \\
    \left\langle
        \delta\bm{M}
    \right\rangle_{\mathrm{obs}}
    &=
    \bm{0}.
\end{aligned}
    \label{eq:mobility_mean_fluctuation_decomposition}
\end{equation}
This separates the mean transport structure, encoded by the constant matrix $\bm{\Phi}$, from deviations around it.

Using the drift fields \(\bm{b}_{\bm{A},n,j}\) defined in
\eqref{eq:b_A_n_j_definition}, the correlation identity can be written
componentwise as the sum of the mean-mobility and correction impulse responses,
\begin{equation}
\begin{aligned}
    \left[\dot{\bm{C}}_{m,n}(t)\right]_{:,j}
    &=
    \left\langle
        \phi_m(\bm{x}(t))\,
        B_{\bm{b}_{\bm{\Phi},n,j}}(\bm{x}(0))
    \right\rangle
    \\
    &\quad+
    \left\langle
        \phi_m(\bm{x}(t))\,
        B_{\bm{b}_{\delta\bm{M},n,j}}(\bm{x}(0))
    \right\rangle .
\end{aligned}
    \label{eq:Cdot_mean_delta_gamma_method}
\end{equation}
This identity holds for each \(j=1,\dots,d_n\).
The first term is the GFDT response associated with the constant mean mobility,
while the second term is the response associated with state-dependent mobility
corrections. The derivation of \eqref{eq:Cdot_mean_delta_gamma_method} is given
in Appendix~\ref{app:mean_mobility_correction_derivation}.

The mean mobility is fixed by the right derivative at the origin of the
coordinate correlation. Taking
\begin{equation}
    \phi_m(\bm{x})=\phi_n(\bm{x})=\bm{x},
    \label{eq:coordinate_observable_choice_method}
\end{equation}
one obtains
\begin{equation}
    \dot{\bm{C}}_{\bm{x},\bm{x}}(0^+)
    =
    -
    \left\langle
        \bm{M}
    \right\rangle_{\mathrm{obs}}
    =
    -
    \bm{\Phi}.
    \label{eq:Phi_from_Cdot_zero_method}
\end{equation}
Thus,
\begin{equation}
    \bm{\Phi}
    =
    -
    \dot{\bm{C}}_{\bm{x},\bm{x},\mathrm{obs}}(0^+).
    \label{eq:Phi_estimator_method}
\end{equation}
In practice, the derivative in \eqref{eq:Phi_estimator_method} is estimated from a smooth short-lag fit of the observed coordinate correlations.

We now keep the coordinate observables and consider finite lags $t>0$. Let
\begin{equation}
    \bm{m}_t(\bm{x}_0)
    :=
    \mathbb{E}
    \left[
        \bm{x}(t)
        \,\middle|\,
        \bm{x}(0)=\bm{x}_0
    \right]
    \label{eq:conditional_mean_method}
\end{equation}
be the finite-lag conditional mean. Then (see Appendix~\ref{app:mean_mobility_correction_derivation})
\begin{equation}
\begin{aligned}
    \dot{\bm{C}}_{\bm{x},\bm{x}}(t)
    &=
    \left\langle
        \bm{x}_t\,
        \bm{s}(\bm{x}_0)^T
    \right\rangle_{\mathrm{obs}}
    \bm{\Phi}
    \\
    &\quad-
    \left\langle
        \nabla\bm{m}_t(\bm{x}_0)\,
        \delta\bm{M}(\bm{x}_0)
    \right\rangle .
\end{aligned}
    \label{eq:coordinate_Cdot_mean_delta_method}
\end{equation}
Here \(\bm{s}\) is the stationary score defined in \eqref{eq:score_definition_method}. The second term is the only contribution of the state-dependent mobility correction to the coordinate correlation dynamics. It vanishes whenever the conditional mean is affine in the initial condition, as occurs, for example, when the joint distribution of $(\bm{x}(t),\bm{x}(0))$ is multivariate Gaussian. In fact
\begin{equation}
    \bm{m}_t(\bm{x}_0)
    =
    \bm{a}(t)
    +
    \bm{A}(t)\bm{x}_0,
    \label{eq:affine_conditional_mean_method}
\end{equation}
then $\nabla\bm{m}_t=\bm{A}(t)$ is independent of $\bm{x}_0$, and therefore
\begin{equation}
    \left\langle
        \nabla\bm{m}_t(\bm{x}_0)\,
        \delta\bm{M}(\bm{x}_0)
    \right\rangle
    =
    \bm{A}(t)
    \left\langle
        \delta\bm{M}
    \right\rangle
    =
    \bm{0}.
    \label{eq:deltaM_vanishes_affine_method}
\end{equation}
Consequently, when the modeling target is to reproduce the observed invariant density and the coordinate time correlations, and when the observed conditional mean is well approximated by an affine function of the initial condition, the constant-mobility model
\begin{equation}
    \bm{M}(\bm{x})
    =
    \bm{\Phi}
    \label{eq:constant_mobility_closure_method}
\end{equation}
is sufficient at the level of these constraints. This recovers, as a special case, the constant score-based Langevin closure introduced in \cite{giorgini2026scorelangevin}. State-dependent corrections $\delta\bm{M}$ become necessary when the selected observables, nonlinear correlations, or non-affine conditional means reveal dynamical features that cannot be represented by the constant closure.

When the constant closure is not sufficiently expressive, we represent the state-dependent correction with a neural parametrization of the full mobility. We write
\begin{equation}
\begin{aligned}
    \bm{M}_\theta(\bm{x})
    &=
    \bm{D}_\theta(\bm{x})
    +
    \bm{R}_\theta(\bm{x}),
    \\
    \bm{D}_\theta(\bm{x})^T
    &=
    \bm{D}_\theta(\bm{x}),
    \\
    \bm{R}_\theta(\bm{x})^T
    &=
    -\bm{R}_\theta(\bm{x}).
\end{aligned}
    \label{eq:neural_M_DR_parametrization}
\end{equation}
A neural network outputs
\begin{equation}
    \bm{h}_\theta(\bm{x})
    =
    \left(
        \bm{\ell}_\theta(\bm{x}),
        \bm{r}_\theta(\bm{x})
    \right)
    \in
    \mathbb{R}^{D(D+1)/2}
    \times
    \mathbb{R}^{D(D-1)/2}.
    \label{eq:neural_output_split}
\end{equation}
The vector $\bm{\ell}_\theta$ fills the entries of a lower-triangular matrix
$\bm{L}_\theta(\bm{x})$, with positive diagonal entries imposed through a smooth positive map. We then set
\begin{equation}
    \bm{D}_\theta(\bm{x})
    =
    \bm{L}_\theta(\bm{x})
    \bm{L}_\theta(\bm{x})^T
    +
    \varepsilon\bm{I},
    \qquad
    \varepsilon>0.
    \label{eq:neural_D_psd_parametrization}
\end{equation}
The vector $\bm{r}_\theta$ fills the strictly lower-triangular entries of
$\bm{R}_\theta$, while the remaining entries are fixed by antisymmetry:
\begin{equation}
    [\bm{R}_\theta(\bm{x})]_{ij}
    =
    -
    [\bm{R}_\theta(\bm{x})]_{ji},
    \qquad
    [\bm{R}_\theta(\bm{x})]_{ii}=0 .
    \label{eq:neural_R_minimal_parametrization}
\end{equation}
This parametrization uses
\begin{equation}
    \frac{D(D+1)}{2}
    +
    \frac{D(D-1)}{2}
    =
    D^2
\end{equation}
outputs, exactly matching the number of degrees of freedom of a general mobility matrix, while guaranteeing
\begin{equation}
    \operatorname{Sym}\bm{M}_\theta(\bm{x})
    =
    \bm{D}_\theta(\bm{x})
    \succ 0
\end{equation}
on the observed support.

The learned correction is defined relative to the mean mobility fixed by the short-time coordinate correlations:
\begin{equation}
    \delta\bm{M}_\theta(\bm{x})
    :=
    \bm{M}_\theta(\bm{x})
    -
    \bm{\Phi}.
    \label{eq:deltaM_theta_definition}
\end{equation}
Using the linearity of the empirical correlation operator, the residual constraint takes the form
\begin{equation}
    \widehat{\mathcal{A}}_{m,n,t}
    [\delta\bm{M}_\theta]
    \approx
    \dot{\bm{C}}_{m,n,\mathrm{obs}}(t)
    -
    \widehat{\mathcal{A}}_{m,n,t}[\bm{\Phi}].
    \label{eq:deltaM_residual_constraint}
\end{equation}
The parameters are determined by minimizing
\begin{equation}
\begin{aligned}
    \mathcal{J}_{\delta M}(\theta)
    &=
    \sum_{t\in\mathcal{T}}
    \sum_{m,n}
    w_{m,n,t}
    \Bigl\|
        \widehat{\mathcal{A}}_{m,n,t}[\delta\bm{M}_\theta]
    \Bigr.
    \\
    &\quad-\Bigl.
        \left[
            \dot{\bm{C}}_{m,n,\mathrm{obs}}(t)
            -
            \widehat{\mathcal{A}}_{m,n,t}[\bm{\Phi}]
        \right]
    \Bigr\|_F^2
    \\
    &\quad+
    \lambda_{\mathrm{mean}}
    \left\|
        \left\langle
            \bm{M}_\theta
        \right\rangle_{\mathrm{obs}}
        -
        \bm{\Phi}
    \right\|_F^2
    \\
    &\quad+
    \lambda_{\mathrm{reg}}\mathcal{R}(\theta).
\end{aligned}
\label{eq:deltaM_neural_loss}
\end{equation}
The mean-matching term enforces
\begin{equation}
    \left\langle
        \delta\bm{M}_\theta
    \right\rangle_{\mathrm{obs}}
    \approx
    \bm{0},
    \label{eq:deltaM_mean_constraint_neural}
\end{equation}
so that the network learns only the state-dependent part not already determined by
$\bm{\Phi}$.

Let $\mathcal{B}_t$ be a minibatch of observed lagged pairs
$(\bm{x}_0^{(k)},\bm{x}_t^{(k)})$. The empirical residual operator is then evaluated as
\begin{equation}
\begin{aligned}
    \widehat{\mathcal{A}}^{\mathcal{B}}_{m,n,t}
    [\delta\bm{M}_\theta]
    &=
    -
    \frac{1}{|\mathcal{B}_t|}
    \sum_{k\in\mathcal{B}_t}
    \phi_m(\bm{x}_t^{(k)})
    \\
    &\quad\times
    \widehat{\bm{s}}_{t|0}
    \bigl(
        \bm{x}_t^{(k)}
        \mid
        \bm{x}_0^{(k)}
    \bigr)^T
    \\
    &\quad\times
    \delta\bm{M}_\theta(\bm{x}_0^{(k)})
    \nabla\phi_n(\bm{x}_0^{(k)})^T .
\end{aligned}
\label{eq:minibatch_deltaM_operator}
\end{equation}
Thus training requires only observed lagged pairs, the estimated conditional scores, the observable gradients, and evaluations of the neural mobility at the initial points. No forward simulation of the learned stochastic model enters the loss.

\section{Results}
\label{sec:results}

\subsection{Analytic warmup: the Cox--Ingersoll--Ross square-root diffusion}
\label{subsec:cir_benchmark}

We begin with a one-dimensional benchmark for which all objects entering the inverse construction, the stationary score, the conditional transition score, and the lagged correlation functions, are available in closed form. This makes it possible to verify, without approximation error, that the correlation constraints identify the state-dependent mobility once the observable library removes the relevant nullspace.

We consider the Cox--Ingersoll--Ross square-root diffusion \cite{CoxIngersollRoss1985,Feller1951}
\begin{equation}
    \d X_t
    =
    \kappa(\theta-X_t)\,\d t
    +
    \sqrt{2\gamma X_t}\,\d W_t,
    \qquad X_t>0,
    \label{eq:cir_sde_results}
\end{equation}
with parameters \(\kappa,\theta,\gamma>0\). We assume the usual Feller condition
\(\kappa\theta\geq \gamma\), so that the origin is not reached when the process is initialized in the interior. The invariant density is Gamma
\begin{equation}
    p_{\mathrm{ss}}(x)
    =
    \frac{
        \left(\kappa/\gamma\right)^{\kappa\theta/\gamma}
    }{
        \Gamma(\kappa\theta/\gamma)
    }
    x^{\kappa\theta/\gamma-1}
    \exp\!\left(-\frac{\kappa x}{\gamma}\right),
    \label{eq:cir_stationary_density_results}
\end{equation}
and therefore
\begin{equation}
    s(x)
    =
    \partial_x\log p_{\mathrm{ss}}(x)
    =
    \frac{\kappa\theta-\gamma}{\gamma x}
    -
    \frac{\kappa}{\gamma}.
    \label{eq:cir_stationary_score_results}
\end{equation}
In the score-based representation, the exact mobility is
\begin{equation}
    M(x)=\gamma x.
    \label{eq:cir_true_mobility_results}
\end{equation}
Thus, with the mean-zero decomposition introduced in
Section~\ref{subsec:mean_mobility_and_corrections},
\begin{equation}
\begin{aligned}
    M(x)
    &=
    \Phi+\delta M(x),
    \\
    \Phi
    &=
    \langle M\rangle_{\mathrm{obs}}=\theta\gamma,
    \\
    \delta M(x)
    &=
    \gamma(x-\theta).
\end{aligned}
    \label{eq:cir_mobility_decomposition_results}
\end{equation}

For a mobility
correction \(\delta\bm{M}=\bm{M}-\bm{\Phi}\), the corresponding residual
equation is
\begin{equation}
\begin{aligned}
    \mathcal{A}_{m,n,t}[\delta\bm{M}]
    &=
    \mathcal{E}_{m,n}(t),
    \\
    \mathcal{E}_{m,n}(t)
    &:=
    \dot{\bm{C}}_{m,n}(t)
    -
    \mathcal{A}_{m,n,t}[\bm{\Phi}].
\end{aligned}
    \label{eq:results_deltaM_residual}
\end{equation}

The coordinate observable \(\phi_1(x)=x\) identifies the constant part immediately. Indeed, using \eqref{eq:Phi_from_Cdot_zero_method},
\begin{equation}
    \Phi
    =
    -\dot C_{1,1,\mathrm{obs}}(0^+)
    =
    \theta\gamma.
    \label{eq:cir_phi_recovery_results}
\end{equation}
However, this observable does not identify the state-dependent correction. If one restricts the correction to the affine mean-zero family
\begin{equation}
    \delta M_a(x)=a(x-\theta),
    \label{eq:cir_affine_correction_family_results}
\end{equation}
then the coordinate channel lies in the nullspace:
\begin{equation}
    \left\langle
        X_t\,s_{t|0}(X_t\mid X_0)\,
        \delta M_a(X_0)
    \right\rangle_{\mathrm{obs}}
    =
    0.
    \label{eq:cir_coordinate_nullspace_results}
\end{equation}
This degeneracy reflects the affine conditional mean of the CIR process and is precisely the kind of nullspace that the observable library must remove.

To recover the mobility correction, we enlarge the observable library by taking
\begin{equation}
    \phi_\alpha(x)=x^\alpha,
    \qquad \alpha\neq 1.
    \label{eq:cir_power_observables_results}
\end{equation}
The exact lagged correlation derivative is
\begin{equation}
\begin{aligned}
    \dot C_{\alpha,1,\mathrm{obs}}(t)
    &=
    -\alpha\theta\gamma
    \left(\frac{\kappa}{\gamma}\right)^{1-\alpha}
    e^{-\kappa t}
    \\
    &\quad\times
    \frac{
        \Gamma(\alpha+\kappa\theta/\gamma)
    }{
        \Gamma(\kappa\theta/\gamma+1)
    }.
\end{aligned}
    \label{eq:cir_exact_correlation_derivative_results}
\end{equation}
Using the residual identity \eqref{eq:results_deltaM_residual}, the inverse problem for the affine correction reduces to
\begin{equation}
    \mathcal{A}_{t,\alpha,1}[\delta M_a]
    =
    \mathcal{E}_{\alpha,1,\mathrm{obs}}(t),
    \label{eq:cir_affine_inverse_problem_results}
\end{equation}
where all quantities are known analytically. In Appendix~\ref{app:cir_benchmark_details} we show that
\begin{equation}
\begin{aligned}
    \mathcal{A}_{t,\alpha,1}[\delta M_a]
    &=
    -a\,K_\alpha(t),
    \\
    \mathcal{E}_{\alpha,1,\mathrm{obs}}(t)
    &=
    -\gamma\,K_\alpha(t),
\end{aligned}
    \label{eq:cir_affine_exact_reduction_results}
\end{equation}
with \(K_\alpha(t)\not\equiv 0\) for \(\alpha\neq 1\). Therefore
\begin{equation}
\begin{aligned}
    a&=\gamma,
    \\
    \delta M(x)&=\gamma(x-\theta),
    \\
    M(x)&=\theta\gamma+\gamma(x-\theta)=\gamma x.
\end{aligned}
    \label{eq:cir_exact_mobility_recovery_results}
\end{equation}

This example isolates the main mechanism of the method. The short-lag coordinate correlation fixes the constant baseline \(\Phi\), while a nonlinear observable removes the nullspace left by the coordinate channel and recovers the full state-dependent mobility exactly.


\subsection{A two-dimensional nonreversible diffusion with affine multiplicative noise}
\label{subsec:affine_multiplicative_2d_results}

We next test the full data-driven pipeline on a two-dimensional nonequilibrium diffusion. In contrast with the CIR benchmark of Section~\ref{subsec:cir_benchmark}, the stationary score, conditional score, and lagged correlation derivatives cannot be obtained analytically. They are all estimated from trajectory data, and the mobility is learned only through the correlation constraints derived in Section~\ref{subsec:mobility_from_correlation_constraints}. The purpose of this example is therefore to assess whether the proposed inverse formulation can recover a nontrivial state-dependent mobility correction from finite-lag statistics, and whether the resulting reduced Langevin model improves dynamical observables without degrading the invariant density with respect to the constant mobility closure baseline.

The reference process is a two-dimensional It\^o diffusion
\begin{equation}
    \d \bm{x}_t
    =
    \bm{f}(\bm{x}_t)\,\d t
    +
    \bm{B}(\bm{x}_t)\,\d \bm{W}_t,
    \qquad
    \bm{x}_t=(x_t,y_t)^T,
    \label{eq:affine_2d_reference_sde_main}
\end{equation}
with a confining gradient drift, a rotational nonequilibrium component, and affine multiplicative noise. Specifically,
\begin{equation}
\begin{aligned}
    \bm{f}(\bm{x})
    &=
    -\nabla U(\bm{x})
    +
    \omega \bm{J}\bm{x},
    \\
    \bm{J}
    &=
    \begin{pmatrix}
        0 & -1\\
        1 & 0
    \end{pmatrix},
    \\
    \omega&=1.1,
\end{aligned}
    \label{eq:affine_2d_drift_main}
\end{equation}
where
\begin{equation}
    U(x,y)
    =
    \frac{1}{4}x^4
    +
    \frac{1}{2}x^2y^2
    +
    \frac{1}{4}y^4
    +
    \frac{1}{2}x^2
    +
    \frac{1}{2}y^2.
    \label{eq:affine_2d_potential_main}
\end{equation}
The noise amplitude is affine in the state,
\begin{equation}
    \bm{B}(x,y)
    =
    \bm{B}_0+x\bm{B}_1+y\bm{B}_2.
    \label{eq:affine_2d_noise_affine_main}
\end{equation}
The experiment reported here uses the amplified affine-noise branch
\(\alpha_{\mathrm{noise}}=3\), so the matrices in
\eqref{eq:affine_2d_noise_affine_main} are
\begin{equation}
\begin{aligned}
    \bm{B}_0
    &=
    \begin{pmatrix}
        0.60 & 0.10\\
        0.08 & 0.55
    \end{pmatrix},
    \\
    \bm{B}_1
    &=
    \begin{pmatrix}
        0.36 & 0.12\\
        0.15 & -0.06
    \end{pmatrix},
    \\
    \bm{B}_2
    &=
    \begin{pmatrix}
        -0.09 & 0.21\\
        0.18 & 0.30
    \end{pmatrix}.
\end{aligned}
    \label{eq:affine_2d_B_coefficients_main}
\end{equation}
In the score-based convention
\begin{equation*}
    \d \bm{x}_t
    =
    \bm{b}(\bm{x}_t)\,\d t
    +
    \sqrt{2}\,\bm{\Sigma}(\bm{x}_t)\,\d\bm{W}_t,
\end{equation*}
the symmetric diffusion tensor of the reference process is
\begin{equation}
    \bm{D}_{\mathrm{ref}}(\bm{x})
    =
    \frac{1}{2}\bm{B}(\bm{x})\bm{B}(\bm{x})^T.
    \label{eq:affine_2d_reference_diffusion_tensor_main}
\end{equation}
For the ex-post comparison in Fig.~\ref{fig:affine_mobility}, we define the
reference score-based mobility by
\begin{equation}
\begin{aligned}
    \bm{M}_{\mathrm{ref}}(\bm{x})
    &=
    \bm{D}_{\mathrm{ref}}(\bm{x})
    +
    \bm{R}_{\mathrm{ref}}(\bm{x}),
    \\
    \bm{R}_{\mathrm{ref}}(\bm{x})
    &=
    r_{\mathrm{ref}}(\bm{x})
    \begin{pmatrix}
        0 & -1\\
        1 & 0
    \end{pmatrix},
\end{aligned}
    \label{eq:affine_2d_reference_mobility_main}
\end{equation}
where the scalar circulation field \(r_{\mathrm{ref}}\) is recovered only for
diagnostic purposes by solving
\begin{equation*}
    \bm{f}(\bm{x})
    =
    \bm{M}_{\mathrm{ref}}(\bm{x})\bm{s}_{\mathrm{ref}}(\bm{x})
    +
    \nabla\cdot\bm{M}_{\mathrm{ref}}(\bm{x})
\end{equation*}
on the evaluation grid using a stationary score estimate from an independent
long reference rollout. This diagnostic field is never used in the mobility
fit.

The reference dataset is generated with time step
\(\Delta t=10^{-3}\) using the Euler--Maruyama scheme. We integrate \(36\) independent trajectories up to
final time \(T=3000\), save every \(10^{-2}\), and discard the first
\(10\%\) of each trajectory as burn-in. The resulting post-burn-in samples are
treated as stationary observations. From these samples and their lagged pairs,
we estimate the stationary score \(\widehat{\bm{s}}_\psi\), the conditional
score \(\widehat{\bm{s}}_{\tau|0}\), and the lagged correlation derivatives
\(\widehat{\dot{\bm{C}}}_{m,n,\mathrm{obs}}(\tau)\). The score estimators
follow the denoising score-matching constructions summarized in
Appendix~\ref{app:score_estimation_dsm}. The experiment-specific details of
the derivative estimator and the training implementation are given in
Appendix~\ref{app:affine_multiplicative_2d_details}.

The constant baseline is fixed from the coordinate correlation derivative,
\begin{equation}
    \bm{\Phi}
    =
    -
    \widehat{\dot{\bm{C}}}_{1,1,\mathrm{obs}}(0^+),
    \label{eq:affine_2d_phi_main}
\end{equation}
and the state-dependent correction is obtained by training a neural mobility
field with the loss \eqref{eq:deltaM_neural_loss}, that is, by fitting the
empirical residual equation
\begin{equation}
\begin{aligned}
    \widehat{\mathcal{A}}_{m,n,\tau}[\delta\bm{M}]
    &=
    \widehat{\bm{\mathcal E}}_{m,n,\mathrm{obs}}(\tau),
    \\
    \widehat{\bm{\mathcal E}}_{m,n,\mathrm{obs}}(\tau)
    &:=
    \widehat{\dot{\bm{C}}}_{m,n,\mathrm{obs}}(\tau)
    -
    \widehat{\mathcal{A}}_{m,n,\tau}[\bm{\Phi}].
\end{aligned}
    \label{eq:affine_2d_residual_identity_main}
\end{equation}
The fitted observable channels use
\begin{equation}
\begin{aligned}
    \phi_n&\in\{x,y\},
    \\
    \phi_m&\in\{x,y,x^2,xy,y^2,x^3,x^2y,xy^2,y^3\},
\end{aligned}
    \label{eq:affine_2d_observable_library_main}
\end{equation}
so that the mobility is constrained by coordinate, quadratic, and cubic probes of the lagged dynamics.
The first task in this benchmark is therefore to infer \(\bm{M}_\theta\) from
observed stationary samples and lagged pairs alone. The second is to integrate
the resulting score-based Langevin model and compare its steady-state density
and lagged correlations with both the observed data and the constant-\(\bm{\Phi}\)
closure, in order to assess how much dynamical fidelity is gained by learning
the state dependence of the mobility.

\begin{figure*}[t]
    \centering
    \includegraphics[width=0.98\textwidth]{fit_M.png}
    \caption{Reference and learned score-based mobility fields for the two-dimensional affine multiplicative-noise benchmark.}
    \label{fig:affine_mobility}
\end{figure*}

Figure~\ref{fig:affine_mobility} shows the resulting mobility reconstruction.
As seen from the four panels, the learned mobility captures the dominant
spatial organization of the four mobility components on the support of the
observed invariant measure. The agreement is better in \(M_{11}\) and
\(M_{12}\), while \(M_{21}\) is the most difficult component. Quantitatively,
the heatmap-on-support relative RMSE for \(\delta\bm{M}\) is \(0.48\), the
pointwise relative RMSE is \(0.49\), and the mean-correction error is
\(0.011\).

\begin{figure*}[p]
    \centering
    \includegraphics[width=0.92\textwidth]{forward_validation_stats.png}
    \caption{Forward validation for the learned reduced Langevin model. The black curve is the reference process, the blue dashed curve is the full learned model with state-dependent mobility, and the red dotted curve is the constant baseline closure.}
    \label{fig:affine_forward_stats}
\end{figure*}

After training, we simulate the learned score-based Langevin model
\begin{equation}
\begin{aligned}
    \d \bm{x}_t
    &=
    \left[
        \bm{M}_\theta(\bm{x}_t)\widehat{\bm{s}}_\psi(\bm{x}_t)
        +
        \nabla\cdot \bm{M}_\theta(\bm{x}_t)
    \right]\d t
    \\
    &\quad+
    \sqrt{2}\,\bm{\Sigma}_\theta(\bm{x}_t)\,\d \bm{W}_t,
    \\
    \bm{\Sigma}_\theta\bm{\Sigma}_\theta^T
    &=
    \bm{D}_\theta.
\end{aligned}
    \label{eq:affine_2d_learned_rom_main}
\end{equation}
The comparison model is the constant-mobility closure
\begin{equation}
\begin{aligned}
    \d \bm{x}_t
    &=
    \left[
        \bm{\Phi}\widehat{\bm{s}}_\psi(\bm{x}_t)
    \right]\d t
    \\
    &\quad+
    \sqrt{2}\,\bm{\Sigma}_\Phi\,\d\bm{W}_t,
    \\
    \bm{\Sigma}_\Phi\bm{\Sigma}_\Phi^T
    &=
    \frac{\bm{\Phi}+\bm{\Phi}^T}{2}.
\end{aligned}
    \label{eq:affine_2d_constant_phi_rom_main}
\end{equation}
Both reduced models are integrated with the same time step
\(\Delta t=10^{-3}\) used for the reference process. We then evaluate them
against the observed dataset and an independent long reference rollout.
Figure~\ref{fig:affine_forward_stats} reports these forward validations. The
top panels compare the stationary density diagnostics. The plotted density
curves for the full state-dependent model show the same excellent agreement
with the observed data as the constant closure. The bottom panels compare the
coordinate autocorrelations and cross-correlations with the long reference
rollout and the observed data. In all of these channels, the state-dependent
model improves on the constant-\(\bm{\Phi}\) closure.

\begin{figure*}[p]
    \centering
    \includegraphics[width=0.95\textwidth]{forward_validation_cphi.png}
    \caption{Lagged observable correlations for the 18 channels used to constrain the mobility fit.}
    \label{fig:affine_forward_cphi}
\end{figure*}

Figure~\ref{fig:affine_forward_cphi} gives the same comparison on the full set
of fitted observable channels. Across these fitted channels, the learned
state-dependent mobility recovers the lagged-correlation curves much more
faithfully than the constant-$\bm{\Phi}$ closure, especially for the
quadratic and cubic probes. This is the observable family that enters the
loss \eqref{eq:deltaM_neural_loss}, so the figure directly tests the
observable-level adequacy discussion of
Section~\ref{subsec:mobility_from_correlation_constraints}. The learned
\(\bm{M}_\theta\) is selected because it makes the empirical
correlation-derivative residuals small on this finite library, not because it
is guaranteed to coincide pointwise with the diagnostic reference mobility.
The fact that the forward Langevin model reproduces the target correlation
curves much more accurately than the constant closure, despite the residual
pointwise error in Fig.~\ref{fig:affine_mobility}, indicates that the mobility
is not uniquely determined by these constraints on the observed support. What
matters for the reduced-modeling objective is that the learned field lies in a
good equivalence class for the chosen observables and lags, so that the induced
Langevin dynamics reproduce the correlation sector used for training.


\subsection{Stochastic complex-amplitude chain with state-dependent mobility}
\label{subsec:stochastic_complex_amplitude_chain}

As a high-dimensional benchmark with controlled oscillatory dynamics, we consider a periodic chain of stochastic complex amplitudes. Complex-amplitude equations of Ginzburg--Landau and Stuart--Landau type arise as normal forms for spatially extended systems near Hopf bifurcation and are standard models for oscillatory media, pattern formation, and nonequilibrium wave dynamics \cite{Stuart1960,NewellWhitehead1969,KuramotoTsuzuki1976,Kuramoto1984,CrossHohenberg1993,AransonKramer2002}. The example below is a finite-dimensional score-form realization of this class: it retains the translation symmetry and underdamped complex-amplitude dynamics of the complex Ginzburg--Landau setting, while prescribing the invariant density and the state-dependent dissipative and reactive mobilities explicitly.

We take \(K=16\) sites on a periodic one-dimensional lattice. At each site,
\[
    \bm{z}_i(t)
    =
    \begin{pmatrix}
        q_i(t)\\
        p_i(t)
    \end{pmatrix}
    \in\mathbb{R}^2,
    \qquad
    i=1,\ldots,K,
\]
or, equivalently, \(a_i(t)=q_i(t)+\mathrm{i}p_i(t)\). The full state is
\begin{equation}
    \bm{z}(t)
    =
    \bigl(\bm{z}_1(t),\ldots,\bm{z}_K(t)\bigr)
    \in\mathbb{R}^{2K},
    \qquad
    2K=32,
    \label{eq:cgl_chain_state_definition}
\end{equation}
with periodic indexing \(\bm{z}_{i+K}=\bm{z}_i\). We define the confining, translation-invariant potential
\begin{equation}
    U(\bm{z})
    =
    \sum_{i=1}^{K}
    \left[
        \frac{\alpha}{2}\|\bm{z}_i\|^2
        +
        \frac{\beta}{4}\|\bm{z}_i\|^4
        +
        \frac{\kappa}{2}\|\bm{z}_{i+1}-\bm{z}_i\|^2
    \right].
    \label{eq:cgl_chain_potential}
\end{equation}
The target stationary density is
\begin{equation}
\begin{aligned}
    p_{\mathrm{ss}}(\bm{z})
    &=
    Z^{-1}\exp[-U(\bm{z})],
    \\
    \bm{s}(\bm{z})
    &=
    \nabla\log p_{\mathrm{ss}}(\bm{z})
    =
    -\nabla U(\bm{z}).
\end{aligned}
    \label{eq:cgl_chain_stationary_density_score}
\end{equation}
The quartic term makes the density normalizable and provides a mild nonlinear amplitude saturation, while the nearest-neighbor term couples the local complex amplitudes and produces spatially coherent modes.

Let
\begin{equation}
    \bm{J}
    =
    \begin{pmatrix}
        0 & -1\\
        1 & 0
    \end{pmatrix}.
    \label{eq:cgl_chain_J_matrix}
\end{equation}
The mobility is chosen to be local, translation equivariant, and state dependent. For each site, set
\begin{align}
    r_i^2(\bm{z})
    &=
    \|\bm{z}_i\|^2,
    \\
    d_i(\bm{z})
    &=
    d_0+d_1 r_i^2(\bm{z}),
    \\
    \omega_i(\bm{z})
    &=
    \omega_0+\omega_1 r_i^2(\bm{z}).
    \label{eq:cgl_chain_local_mobility_coefficients}
\end{align}
The symmetric and antisymmetric parts of the mobility are
\begin{equation}
    \bm{D}(\bm{z})
    =
    \operatorname{blockdiag}_{i=1}^{K}
    \left[
        d_i(\bm{z})\,\bm{I}_2
    \right],
    \label{eq:cgl_chain_D_definition}
\end{equation}
and
\begin{equation}
    \bm{R}(\bm{z})
    =
    \operatorname{blockdiag}_{i=1}^{K}
    \left[
        \omega_i(\bm{z})\,\bm{J}
    \right].
    \label{eq:cgl_chain_R_definition}
\end{equation}
Thus
\begin{equation}
\begin{aligned}
    \bm{M}(\bm{z})
    &=
    \bm{D}(\bm{z})
    +
    \bm{R}(\bm{z}),
    \\
    \bm{D}(\bm{z})^T=\bm{D}(\bm{z})\succ0,
    &\qquad
    \bm{R}(\bm{z})^T=-\bm{R}(\bm{z}).
\end{aligned}
    \label{eq:cgl_chain_M_definition}
\end{equation}
The scalar field \(d_i\) controls the local dissipative mobility and the noise amplitude, while \(\omega_i\) controls the local reactive rotation. Both depend on the instantaneous amplitude \(r_i^2\), so the true diffusion tensor and the true irreversible transport tensor are state dependent.

The reference process is the score-based It\^o diffusion
\begin{align}
    \d\bm{z}(t)
    &=
    \left[
        \bm{M}(\bm{z}(t))\bm{s}(\bm{z}(t))
        +
        \nabla\cdot\bm{M}(\bm{z}(t))
    \right]\d t
    \\
    &\quad+
    \sqrt{2}\,
    \bm{\Sigma}(\bm{z}(t))\,\d\bm{W}_t,
    \label{eq:cgl_chain_score_form_sde}
\end{align}
where
\begin{equation}
    \bm{\Sigma}(\bm{z})\bm{\Sigma}(\bm{z})^T
    =
    \bm{D}(\bm{z}).
    \label{eq:cgl_chain_noise_amplitude}
\end{equation}
Equivalently, in site variables,
\begin{align}
    \d\bm{z}_i(t)
    &=
    \bm{b}_i(\bm{z}(t))\,\d t
    +
    \sqrt{2d_i(\bm{z}(t))}\,\d\bm{W}_{i,t},
    \\
    &\qquad
    i=1,\ldots,K,
    \label{eq:cgl_chain_site_sde}
\end{align}
where the \(\bm{W}_{i,t}\) are independent two-dimensional standard Wiener processes. The local gradient of the potential is
\begin{align}
    \bm{g}_i(\bm{z})
    &:=
    \nabla_{\bm{z}_i}U(\bm{z})
    \\
    &=
    \alpha\bm{z}_i
    +
    \beta\|\bm{z}_i\|^2\bm{z}_i
    \\
    &\quad+
    \kappa
    \left(
        2\bm{z}_i-\bm{z}_{i-1}-\bm{z}_{i+1}
    \right),
    \label{eq:cgl_chain_potential_gradient}
\end{align}
and the drift is
\begin{equation}
    \bm{b}_i(\bm{z})
    =
    -d_i(\bm{z})\,\bm{g}_i(\bm{z})
    -
    \omega_i(\bm{z})\,\bm{J}\bm{g}_i(\bm{z})
    +
    2d_1\bm{z}_i
    +
    2\omega_1\bm{J}\bm{z}_i.
    \label{eq:cgl_chain_site_drift}
\end{equation}
The first term in \eqref{eq:cgl_chain_site_drift} is the dissipative, gradient-driven component. The second term is the nonconservative reactive component that rotates the local force by \(\pi/2\). The last two terms are the It\^o divergence corrections associated with the state dependence of \(\bm{D}\) and \(\bm{R}\). Since the drift is written in the form \eqref{eq:cgl_chain_score_form_sde}, the density \eqref{eq:cgl_chain_stationary_density_score} is invariant by construction. The antisymmetric part \(\bm{R}\) generates a stationary probability current, so the process is nonreversible unless \(\omega_0=\omega_1=0\).

The model is exactly equivariant under cyclic translations. If \(\bm{T}\) denotes the shift
\begin{equation}
    (\bm{T}\bm{z})_i
    =
    \bm{z}_{i-1},
    \label{eq:cgl_chain_shift_operator}
\end{equation}
then
\begin{align}
    U(\bm{T}\bm{z})
    &=
    U(\bm{z}),
    \\
    \bm{D}(\bm{T}\bm{z})
    &=
    \bm{T}\bm{D}(\bm{z})\bm{T}^T,
    \\
    \bm{R}(\bm{T}\bm{z})
    &=
    \bm{T}\bm{R}(\bm{z})\bm{T}^T.
    \label{eq:cgl_chain_translation_equivariance}
\end{align}
Consequently, the stochastic dynamics are equivariant in law under translations of the periodic lattice. This symmetry makes the process well matched to a one-dimensional periodic U-Net architecture with two channels per lattice site for the stationary score, and with the corresponding pair of initial and final states plus the lag variable for the conditional score.

The parameters used in the experiment are
\begin{equation}
\begin{aligned}
    K &= 16,
    &
    \alpha &= 0.70,
    &
    \beta &= 0.08,
    &
    \kappa &= 0.20,
    \\
    d_0 &= 0.30,
    &
    d_1 &= 0.04,
    &
    \omega_0 &= 2.20,
    &
    \omega_1 &= 0.15.
\end{aligned}
    \label{eq:cgl_chain_parameter_values}
\end{equation}
These values place the deterministic skeleton in a stable underdamped regime while keeping the nonlinear and state-dependent effects visible on the typical support of the invariant density. To see this, linearize around the homogeneous state \(\bm{z}=0\) and decompose the chain into lattice Fourier modes. The Hessian eigenvalue of the potential for mode \(k\) is
\begin{equation}
    \lambda_k
    =
    \alpha
    +
    4\kappa
    \sin^2\left(\frac{\pi k}{K}\right),
    \qquad
    k=0,\ldots,K-1.
    \label{eq:cgl_chain_linear_potential_spectrum}
\end{equation}
The corresponding linearized deterministic dynamics have the form
\begin{equation}
    \frac{\d}{\d t}\widehat{\bm{z}}_k
    =
    -\Gamma_k\widehat{\bm{z}}_k
    -
    \Omega_k\bm{J}\widehat{\bm{z}}_k,
    \label{eq:cgl_chain_linearized_modes}
\end{equation}
with
\begin{equation}
    \Gamma_k
    =
    d_0\lambda_k-2d_1,
    \qquad
    \Omega_k
    =
    \omega_0\lambda_k-2\omega_1.
    \label{eq:cgl_chain_linear_damping_frequency}
\end{equation}
For the parameter values in \eqref{eq:cgl_chain_parameter_values}, one has
\[
    \Gamma_k\in[0.13,0.37],
    \qquad
    \Omega_k\in[1.24,3.00].
\]
Thus all Fourier modes are linearly stable and underdamped. Their coordinate correlations have the leading behavior
\begin{equation}
    C_k(t)
    \sim
    e^{-\Gamma_k t}\cos(\Omega_k t),
    \label{eq:cgl_chain_linear_acf}
\end{equation}
with sine-type cross-correlations between the \(q\) and \(p\) components. The process therefore provides a high-dimensional, translation-invariant, nonchaotic nonequilibrium diffusion whose finite-lag correlations contain coherent decaying oscillations and whose true mobility has both state-dependent symmetric and antisymmetric parts.


\subsection{A periodic fluctuating-hydrodynamic chain with acoustic correlations}
\label{subsec:fluctuating_hydrodynamic_chain}

As a high-dimensional non-chaotic benchmark with structured
state-dependent transport, we consider a finite-volume
discretization of the
one-dimensional isothermal Landau--Lifshitz fluctuating hydrodynamics on a
periodic domain. Fluctuating hydrodynamics augments the deterministic
Navier--Stokes equations by stochastic fluxes whose covariance is fixed by a
fluctuation--dissipation relation \cite{LandauLifshitzFluidMechanics,
BellGarciaWilliams2007,DonevVandenEijndenGarciaBell2010}. This model is a
standard mesoscopic description of thermal fluctuations in fluids, and its
finite-volume discretizations have been studied extensively in the numerical
fluctuating-hydrodynamics literature. It is also naturally compatible with the
reversible--irreversible decomposition used throughout this paper: the inviscid
Euler part is generated by a state-dependent antisymmetric Poisson operator,
while viscosity and stochastic stress are generated by a symmetric positive
semidefinite friction operator, as in the GENERIC formulation of nonequilibrium
thermodynamics \cite{GrmelaOttinger1997,OttingerGrmela1997,Ottinger2005}.

We use a one-dimensional periodic box of length \(L=2\pi\), divided into
\(N=8\) finite-volume cells of width
\begin{equation}
    \Delta x = \frac{L}{N}.
    \label{eq:fhd_dx}
\end{equation}
The cell variables are the mass density \(\rho_i(t)>0\) and momentum density
\(m_i(t)\), with velocity
\begin{equation}
    u_i(t)=\frac{m_i(t)}{\rho_i(t)},
    \qquad
    i=1,\ldots,N.
    \label{eq:fhd_velocity_definition}
\end{equation}
The full conservative state is
\begin{equation}
    \bm{y}(t)
    =
    \bigl(\rho_1(t),\ldots,\rho_N(t),
          m_1(t),\ldots,m_N(t)\bigr)^T
    \in \mathbb{R}^{2N}.
    \label{eq:fhd_state_definition}
\end{equation}
Periodic indexing is used throughout. The total mass and total momentum are
conserved by the periodic fluctuating-hydrodynamic dynamics. We therefore fix
\begin{equation}
    \Delta x\sum_{i=1}^{N}\rho_i = \rho_0 L,
    \qquad
    \Delta x\sum_{i=1}^{N}m_i = 0,
    \label{eq:fhd_conserved_quantities}
\end{equation}
and regard the process as evolving on this affine subspace. Equivalently, the
two zero Fourier modes are projected out when estimating scores and correlation
statistics. With \(N=8\), the effective dimension is \(2N-2=14\), while the data
retain the periodic lattice structure needed by a one-dimensional translation
equivariant score network.

The equilibrium free energy is taken to be the isothermal ideal-fluid free
energy plus kinetic energy,
\begin{equation}
    \mathcal{F}(\bm{y})
    =
    \Delta x
    \sum_{i=1}^{N}
    \left[
        c_s^2\!
        \left\{
            \rho_i\log\left(\frac{\rho_i}{\rho_0}\right)
            -\rho_i+\rho_0
        \right\}
        +
        \frac{m_i^2}{2\rho_i}
    \right],
    \label{eq:fhd_free_energy}
\end{equation}
where \(c_s\) is the isothermal sound speed and \(\rho_0\) is the reference
density. The corresponding equilibrium density on the constrained subspace is
\begin{equation}
    p_{\mathrm{ss}}(\bm{y})
    =
    Z^{-1}
    \exp\left[-\frac{\mathcal{F}(\bm{y})}{\Theta}\right],
    \label{eq:fhd_stationary_density}
\end{equation}
where \(\Theta\) is the dimensionless thermal energy. Thus the stationary score
is
\begin{align}
    \bm{s}(\bm{y})
    &=
    \nabla_{\bm{y}}\log p_{\mathrm{ss}}(\bm{y})
    \\
    \bm{s}(\bm{y})
    &=
    -\frac{1}{\Theta}\nabla_{\bm{y}}\mathcal{F}(\bm{y}).
    \label{eq:fhd_stationary_score}
\end{align}
with the gradient understood tangentially to the constrained subspace
\eqref{eq:fhd_conserved_quantities}.

We now define the finite-volume operators. For a cell field \(a_i\), let
\begin{equation}
    \bar a_{i+1/2}
    =
    \frac{a_i+a_{i+1}}{2}
    \label{eq:fhd_edge_average}
\end{equation}
denote the centered edge average, and for an edge field \(g_{i+1/2}\), let
\begin{equation}
    (\nabla_h\cdot g)_i
    =
    \frac{g_{i+1/2}-g_{i-1/2}}{\Delta x}
    \label{eq:fhd_discrete_divergence}
\end{equation}
denote the periodic finite-volume divergence. The isothermal pressure is
\begin{equation}
    p(\rho)=c_s^2\rho.
    \label{eq:fhd_pressure}
\end{equation}
The reversible mass and momentum fluxes are
\begin{align}
    j_{i+1/2}
    &=
    \bar m_{i+1/2},
    \\
    \Pi^{\mathrm{rev}}_{i+1/2}
    &=
    \bar m_{i+1/2}\bar u_{i+1/2}
    +
    p(\bar\rho_{i+1/2}),
    \label{eq:fhd_reversible_fluxes}
\end{align}
The viscous stress is
\begin{equation}
    \tau_{i+1/2}
    =
    \eta_{i+1/2}
    \frac{u_{i+1}-u_i}{\Delta x},
    \label{eq:fhd_viscous_stress}
\end{equation}
with density-dependent viscosity
\begin{equation}
    \eta_{i+1/2}
    =
    \eta_0
    \left(
        \frac{\bar\rho_{i+1/2}}{\rho_0}
    \right)^{\zeta}.
    \label{eq:fhd_density_dependent_viscosity}
\end{equation}
This density dependence makes the symmetric diffusion tensor genuinely
state dependent, while preserving positivity of the local viscosity.

The semidiscrete fluctuating-hydrodynamic equations are
\begin{align}
    \d \rho_i
    &=
    -(\nabla_h\cdot j)_i\,\d t,
    \label{eq:fhd_density_equation}
\end{align}
and
\begin{align}
    \d m_i
    =
    -(\nabla_h\cdot \Pi^{\mathrm{rev}})_i\,\d t
    \\
    &+
    (\nabla_h\cdot \tau)_i\,\d t
    \\
    &+
    \frac{\sqrt{2\Theta}}{\Delta x}
    \sqrt{\frac{\eta_{i+1/2}}{\Delta x}}
    \d B_{i+1/2}(t)
    \\
    &-
    \frac{\sqrt{2\Theta}}{\Delta x}
    \sqrt{\frac{\eta_{i-1/2}}{\Delta x}}
    \d B_{i-1/2}(t),
    \label{eq:fhd_momentum_equation}
\end{align}
where the \(B_{i+1/2}\) are independent scalar Brownian motions associated with
the stochastic stresses at cell faces. The same viscosity field
\(\eta_{i+1/2}\) appears in the deterministic viscous flux and in the stochastic
stress covariance, giving the finite-dimensional fluctuation--dissipation
balance.

Equivalently, the process can be written in the score-based form
\begin{align}
    \d\bm{y}
    &=
    \left[
        \bm{M}(\bm{y})\bm{s}(\bm{y})
        +
        \nabla\cdot\bm{M}(\bm{y})
    \right]\d t
    \\
    &+
    \sqrt{2}\,
    \bm{\Sigma}(\bm{y})\,\d\bm{W}_t,
    \label{eq:fhd_score_form_sde}
\end{align}
\begin{align}
    \bm{\Sigma}(\bm{y})\bm{\Sigma}(\bm{y})^T
    &=
    \bm{D}(\bm{y}),
    \\
    \bm{M}(\bm{y})
    &=
    \bm{D}(\bm{y})+\bm{R}(\bm{y}),
    \\
    \bm{D}(\bm{y})^T
    &=
    \bm{D}(\bm{y})\succeq0,
    \\
    \bm{R}(\bm{y})^T
    &=
    -\bm{R}(\bm{y}).
    \label{eq:fhd_mobility_decomposition}
\end{align}
Let \(\bm{B}_h\) denote the periodic edge-to-cell divergence matrix associated
with \eqref{eq:fhd_discrete_divergence}. In conservative variables, the
symmetric part has the block structure
\begin{equation}
    \bm{D}(\bm{y})
    =
    \Theta
    \begin{pmatrix}
        \bm{0} & \bm{0}\\
        \bm{0} &
        \bm{B}_h
        \operatorname{diag}
        \left(
            \frac{\eta_{1/2}}{\Delta x},
            \ldots,
            \frac{\eta_{N+1/2}}{\Delta x}
        \right)
        \bm{B}_h^T
    \end{pmatrix}.
    \label{eq:fhd_symmetric_mobility}
\end{equation}
This matrix is positive semidefinite and state dependent through
\eqref{eq:fhd_density_dependent_viscosity}. The antisymmetric part is
\begin{equation}
    \bm{R}(\bm{y})
    =
    -\Theta\,\bm{L}_h(\bm{y}),
    \label{eq:fhd_antisymmetric_mobility}
\end{equation}
where \(\bm{L}_h(\bm{y})\) is the centered finite-volume discretization of the
noncanonical Euler Poisson operator. Its action on
\(\nabla_{\bm{y}}\mathcal{F}\) gives precisely the reversible flux terms in
\eqref{eq:fhd_density_equation}--\eqref{eq:fhd_momentum_equation},
\begin{equation}
    \bm{R}(\bm{y})\bm{s}(\bm{y})
    =
    \begin{pmatrix}
        -\nabla_h\cdot j\\
        -\nabla_h\cdot\Pi^{\mathrm{rev}}
    \end{pmatrix}.
    \label{eq:fhd_reversible_score_action}
\end{equation}
Thus \(\bm{R}\) is the finite-dimensional reversible transport tensor, while
\(\bm{D}\) is the finite-dimensional dissipative fluctuation tensor. Both are
translation equivariant and local on the periodic lattice; \(\bm{R}\) is
state dependent through the density and momentum dependence of the Euler fluxes,
and \(\bm{D}\) is state dependent through the density-dependent viscosity.

The parameter values used in the experiment are
\begin{align}
    N &= 8,
    \quad
    L = 2\pi,
    \quad
    \rho_0 = 1,
    \\
    c_s &= 1.2,
    \quad
    \Theta &= 2.5\times 10^{-2},
    \\
    \eta_0 &= 8.0\times 10^{-2},
    \quad
    \zeta &= \frac{1}{2}.
    \label{eq:fhd_parameter_values}
\end{align}
These values keep the process in a weakly compressible, single-phase regime,
so that the density remains close to \(\rho_0\), while preserving visible
state dependence in the viscous and stochastic stress amplitudes.

The reason this benchmark is dynamically well suited to the present method is
that the deterministic dynamics are acoustic rather than chaotic. Linearizing
\eqref{eq:fhd_density_equation}--\eqref{eq:fhd_momentum_equation} about the
homogeneous equilibrium \(\rho_i=\rho_0\), \(m_i=0\), each nonzero Fourier mode
satisfies, to leading order,
\begin{align}
    \frac{\d}{\d t}
    \begin{pmatrix}
        \delta\widehat{\rho}_k\\
        \widehat{m}_k
    \end{pmatrix}
    =
    \begin{pmatrix}
        0 & -\mathrm{i}k_h\\
        -\mathrm{i}c_s^2 k_h & -\nu k_h^2
    \end{pmatrix}
    \begin{pmatrix}
        \delta\widehat{\rho}_k\\
        \widehat{m}_k
    \end{pmatrix}
    ,
    \\
    \nu=\frac{\eta_0}{\rho_0},
    \label{eq:fhd_linearized_mode_equation}
\end{align}
where
\begin{equation}
    k_h
    =
    \frac{2}{\Delta x}
    \sin\left(\frac{\pi k}{N}\right),
    \qquad
    k=1,\ldots,N-1.
    \label{eq:fhd_discrete_wavenumber}
\end{equation}
The corresponding eigenvalues are
\begin{equation}
    \lambda_k^{\pm}
    =
    -\frac{\nu k_h^2}{2}
    \pm
    \mathrm{i}
    \left[
        c_s^2 k_h^2
        -
        \frac{\nu^2 k_h^4}{4}
    \right]^{1/2}.
    \label{eq:fhd_acoustic_eigenvalues}
\end{equation}
For the parameters in \eqref{eq:fhd_parameter_values}, the nonzero modes are
underdamped. Consequently, density and momentum correlations contain coherent
decaying oscillations associated with sound propagation and viscous attenuation,
rather than the rapid chaotic decorrelation present in strongly nonlinear
lattice systems. This provides a translation-invariant, physically motivated,
high-dimensional nonequilibrium diffusion with state-dependent symmetric and
antisymmetric mobility tensors.

% \subsection{A high-dimensional stochastic Lorenz--96 system with translation-invariant correlated noise}
% \label{subsec:l96_correlated_noise}

% As a high-dimensional test case, we consider a stochastic version of the Lorenz--96 model. The Lorenz--96 system is a canonical finite-dimensional model for atmospheric predictability, data assimilation, and reduced stochastic modeling \cite{Lorenz1996,LorenzEmanuel1998}. Although idealized, it retains several structural features that are central in geophysical turbulence: nonlinear advection, linear damping, constant large-scale forcing, chaotic energy transfer across spatial scales, and translation symmetry on a periodic one-dimensional lattice. Stochastic versions of the model are also widely used as test beds for parameterizing unresolved degrees of freedom and model error \cite{MajdaTimofeyevVandenEijnden1999,Wilks2005,Palmer2009,Berner2017}.

% We take \(K=40\) resolved variables,
% \begin{equation}
%     \bm{x}(t)
%     =
%     \bigl(x_1(t),\ldots,x_K(t)\bigr)^T
%     \in \mathbb{R}^K,
%     \qquad
%     K=40,
%     \label{eq:l96_state_definition}
% \end{equation}
% with periodic indexing,
% \begin{equation}
%     x_{i+K}=x_i,
%     \qquad
%     i\in\mathbb{Z}.
%     \label{eq:l96_periodic_indexing}
% \end{equation}
% The deterministic Lorenz--96 vector field is
% \begin{equation}
%     f_i(\bm{x})
%     =
%     x_{i-1}
%     \bigl(x_{i+1}-x_{i-2}\bigr)
%     -
%     x_i
%     +
%     F,
%     \qquad
%     i=1,\ldots,K,
%     \label{eq:l96_deterministic_vector_field}
% \end{equation}
% and we use the standard chaotic forcing value
% \begin{equation}
%     F=8.
%     \label{eq:l96_forcing_value}
% \end{equation}
% The stochastic model is then
% \begin{equation}
%     \d x_i(t)
%     =
%     f_i(\bm{x}(t))\,\d t
%     +
%     \sqrt{2}
%     \sum_{j=1}^{K}
%     \bm{Q}^{1/2}_{ij}\,
%     \d W_j(t),
%     \qquad
%     i=1,\ldots,K,
%     \label{eq:l96_stochastic_component_form}
% \end{equation}
% or, equivalently,
% \begin{equation}
%     \d \bm{x}(t)
%     =
%     \bm{f}(\bm{x}(t))\,\d t
%     +
%     \sqrt{2}\,
%     \bm{Q}^{1/2}\,\d \bm{W}_t,
%     \label{eq:l96_stochastic_vector_form}
% \end{equation}
% where \(\bm{W}_t\) is a standard \(K\)-dimensional Wiener process. In the notation of Section~\ref{sec:method}, the diffusion tensor is constant and given by
% \begin{equation}
%     \bm{D}(\bm{x})
%     =
%     \bm{Q}.
%     \label{eq:l96_diffusion_tensor}
% \end{equation}
% Thus this example tests the method in a genuinely high-dimensional setting while retaining an additive, analytically controlled diffusion structure.

% The noise covariance is chosen to be spatially homogeneous and correlated on the periodic lattice. Let \(\bm{U}\) denote the unitary discrete Fourier transform on \(K\) grid points, and let
% \begin{equation}
%     \lambda_m
%     =
%     4\sin^2\left(\frac{\pi m}{K}\right),
%     \qquad
%     m=0,\ldots,K-1,
%     \label{eq:l96_laplacian_eigenvalues}
% \end{equation}
% be the eigenvalues of the negative periodic discrete Laplacian. We define
% \begin{equation}
%     \bm{Q}
%     =
%     \bm{U}^*
%     \operatorname{diag}
%     \bigl(q_0,\ldots,q_{K-1}\bigr)
%     \bm{U},
%     \label{eq:l96_noise_covariance_fourier}
% \end{equation}
% with spectrum
% \begin{equation}
%     q_m
%     =
%     \sigma^2
%     \left[
%         (1-\rho)
%         +
%         \rho\,
%         \frac{
%             (1+\ell^2\lambda_m)^{-p}
%         }{
%             \frac{1}{K}
%             \sum_{n=0}^{K-1}
%             (1+\ell^2\lambda_n)^{-p}
%         }
%     \right].
%     \label{eq:l96_noise_spectrum}
% \end{equation}
% The parameters used in the experiment are
% \begin{equation}
%     \sigma=0.5,
%     \qquad
%     \rho=0.6,
%     \qquad
%     \ell=2,
%     \qquad
%     p=2.
%     \label{eq:l96_noise_parameters}
% \end{equation}
% The normalization in \eqref{eq:l96_noise_spectrum} gives
% \begin{equation}
%     \frac{1}{K}\operatorname{tr}\bm{Q}
%     =
%     \frac{1}{K}
%     \sum_{m=0}^{K-1}q_m
%     =
%     \sigma^2,
%     \label{eq:l96_noise_trace_normalization}
% \end{equation}
% so that \(\sigma\) fixes the root-mean-square diffusion amplitude per resolved variable, independently of the correlation parameters \(\rho\), \(\ell\), and \(p\). The parameter \(\rho\) controls the fraction of the variance placed in the spatially correlated component, \(\ell\) sets the correlation length on the periodic lattice, and \(p\) controls the spectral decay. The nonzero spectral floor \(1-\rho\) prevents the forcing from becoming rank deficient and ensures that all Fourier modes remain directly forced.

% This choice of \(\bm{Q}\) is designed to respect the symmetries of the Lorenz--96 dynamics. Let \(\bm{T}\) be the cyclic shift operator,
% \begin{equation}
%     (\bm{T}\bm{x})_i
%     =
%     x_{i-1}.
%     \label{eq:l96_shift_operator}
% \end{equation}
% The deterministic vector field satisfies
% \begin{equation}
%     \bm{f}(\bm{T}\bm{x})
%     =
%     \bm{T}\bm{f}(\bm{x}),
%     \label{eq:l96_drift_equivariance}
% \end{equation}
% and the circulant covariance \eqref{eq:l96_noise_covariance_fourier} satisfies
% \begin{equation}
%     \bm{T}\bm{Q}\bm{T}^T
%     =
%     \bm{Q}.
%     \label{eq:l96_covariance_translation_invariance}
% \end{equation}
% Consequently, the stochastic dynamics are equivariant in law under cyclic translations of the lattice. No grid point is distinguished by either the deterministic dynamics or the stochastic forcing.

% The physical interpretation is that the stochastic term represents the net effect of unresolved processes or model error on the resolved Lorenz--96 variables. Independent noise at each lattice site would correspond to unresolved forcing with zero spatial coherence. By contrast, the covariance \eqref{eq:l96_noise_covariance_fourier} represents unresolved fluctuations with a finite correlation length, as expected when unresolved degrees of freedom influence neighboring resolved variables coherently. The low-pass spectral envelope models the fact that unresolved feedbacks often project more strongly onto large and intermediate resolved scales than onto grid-scale oscillations, while the isotropic spectral floor represents residual small-scale model error and guarantees nondegenerate stochastic excitation. The resulting process is therefore a translation-invariant, high-dimensional nonequilibrium diffusion with physically interpretable correlated additive noise.

\section{Conclusions}
\label{sec:conclusions}

We have introduced a conditional score-based framework for constructing
effective Langevin models from stationary samples and finite-lag observations.
The approach separates the modeling problem into two complementary components:
the stationary score, which encodes the invariant density, and a mobility field,
which determines the dynamical transport structure of the reduced stochastic
model. Once inferred, the mobility specifies both the diffusion tensor and the
drift term required to preserve the observed invariant density, including the
irreversible probability current in nonequilibrium systems.

The main result of the paper is the inverse relation in
Eq.~\eqref{eq:central_conditional_score_identity}, which constrains the mobility
directly from lagged correlation derivatives and conditional scores. The same
identity is also a finite-lag GFDT statement: it identifies the selected
correlation derivative with the impulse response to the drift direction
\(-\bm{M}\nabla\phi_n^T\). The relation is linear in the unknown mobility and
therefore provides the basis for a direct calibration procedure. In contrast
with simulation-based calibration,
the method does not require repeated forward integrations of candidate reduced
models. In contrast with local Kramers--Moyal or finite-difference
reconstruction methods, it does not require short-time derivative estimates or
finely resolved trajectories. Instead, it uses finite-lag statistical
information, which can be extracted from stationary lagged pairs and is
therefore naturally suited to data that are temporally coarse, unevenly sampled,
or noisy.

The numerical examples illustrate the scope and interpretation of the method.
In the CIR benchmark of
Section~\ref{subsec:cir_benchmark}, the framework shows which mobility components
are constrained by the available correlation identities.
In the affine nonequilibrium example
Section~\ref{subsec:affine_multiplicative_2d_results}, where stationary and
conditional scores are estimated from data, the learned state-dependent mobility
improves finite-lag statistics relative to simpler constant-mobility closures,
while preserving the target invariant distribution.
These results support the central modeling principle of the paper: for reduced
stochastic modeling, the goal is not necessarily to identify a unique pointwise
mobility throughout state space, but rather to learn a transport structure that
reproduces the relevant stationary and finite-time statistics on the portion of
state space sampled by the data.

A broader implication of the work is the role of the conditional score as a
practical representation of finite-time dynamical information. The central
identity of the paper can be interpreted as a Koopman-gradient relation: it
expresses the gradient of an observable evolved under the Koopman semigroup
through averages over lagged pairs. In traditional Koopman spectral approaches,
such quantities are often accessed by constructing eigenfunction expansions of
the Koopman operator. While this perspective is powerful and interpretable, it
becomes difficult to use as a scalable computational tool in high-dimensional
stochastic systems, where many modes may be required and differentiating learned
eigenfunctions can be unstable. The conditional-score formulation provides an
alternative route: it evaluates the same response-relevant Koopman-gradient
information directly from transition statistics, without requiring an explicit
spectral decomposition.

Conditional scores provide a way to rewrite finite-lag averages in terms of
spatial derivatives of transition statistics, without explicitly reconstructing
the transition density or a spectral representation of the Koopman operator. In
this sense, the framework developed here connects score-based generative
modeling, stochastic reduced-order modeling, and response-theoretic Koopman
descriptions of dynamical systems. This connection opens a path toward
data-driven stochastic closures that simultaneously preserve invariant
statistics, reproduce finite-lag correlations, and encode response-relevant
transport information in a form that can scale to high-dimensional systems.

Several directions remain open. The present formulation can be extended by using
richer parameterizations of the mobility, by incorporating regularization that
enforces physical structure such as symmetry, positive semidefiniteness,
sparsity, or locality, and by adapting conditional score estimation to strongly
multiscale or partially observed systems. Another important direction is to
develop systematic criteria for choosing observable libraries and lag times so
that the inferred mobility is well constrained for the dynamical features of
interest. Finally, the Koopman interpretation suggests a broader program in
which conditional score networks are used as differential surrogates for
finite-time Koopman or Markov semigroups, providing a scalable alternative to
explicit eigenfunction-based constructions in response theory and stochastic
model reduction.

\section*{Acknowledgments}

The author thanks Andre N. Souza, Jonah Botvinick-Greenhouse, and Yuanchao Xu for useful discussions.

\clearpage
\onecolumn
\newgeometry{margin=1in}
\appendix
\setcounter{section}{0}
\setcounter{subsection}{0}
\setcounter{equation}{0}
\setcounter{figure}{0}
\setcounter{table}{0}
\renewcommand{\thesection}{S\arabic{section}}
\renewcommand{\thesubsection}{S\arabic{section}.\arabic{subsection}}
\renewcommand{\theequation}{S\arabic{equation}}
\renewcommand{\thefigure}{S\arabic{figure}}
\renewcommand{\thetable}{S\arabic{table}}
\renewcommand{\theHsection}{appendix.\arabic{section}}
\renewcommand{\theHsubsection}{appendix.\arabic{section}.\arabic{subsection}}
\renewcommand{\theHequation}{appendix.\arabic{equation}}
\renewcommand{\theHfigure}{appendix.\arabic{figure}}
\renewcommand{\theHtable}{appendix.\arabic{table}}
\section*{Appendix}
\section{Score-based representation of stationary diffusions}
\label{app:mobility_representation}

In this appendix we derive the score-based form of the effective diffusion used in the main text. We work on a domain $\Omega\subseteq\mathbb{R}^D$, either $\Omega=\mathbb{R}^D$ with sufficient decay at infinity or a domain with boundary conditions such that all boundary fluxes vanish. We assume that the stationary density $p_{\mathrm{ss}}$ is strictly positive and sufficiently smooth on $\Omega$, and that the drift and diffusion coefficients below are smooth enough for the integrations by parts and differentiations used here to be justified.

Consider the It\^o diffusion
\begin{equation}
    \d \bm{x}(t)
    =
    \bm{F}(\bm{x}(t))\,\d t
    +
    \sqrt{2}\,\bm{\Sigma}(\bm{x}(t))\,\d \bm{W}_t ,
    \label{eq:app_general_sde}
\end{equation}
where $\bm{W}_t$ is a standard $D$-dimensional Wiener process. We define the diffusion tensor
\begin{equation}
    \bm{D}(\bm{x})
    :=
    \bm{\Sigma}(\bm{x})\bm{\Sigma}(\bm{x})^T .
    \label{eq:app_diffusion_tensor}
\end{equation}
Because of the convention $\sqrt{2}\,\bm{\Sigma}$ in \eqref{eq:app_general_sde}, the infinitesimal generator acting on smooth test functions $f$ is
\begin{equation}
    \mathcal{L}f
    =
    \bm{F}\cdot\nabla f
    +
    \bm{D}:\nabla\nabla f
    =
    F_i\,\partial_i f
    +
    D_{ij}\,\partial_i\partial_j f ,
    \label{eq:app_generator}
\end{equation}
where repeated indices are summed over. The corresponding Fokker--Planck equation is
\begin{equation}
    \partial_t p
    =
    \mathcal{L}^*p
    =
    -\partial_i(F_i p)
    +
    \partial_i\partial_j(D_{ij}p).
    \label{eq:app_fokker_planck}
\end{equation}
Equivalently, it can be written as a conservation law
\begin{equation}
    \partial_t p
    =
    -\nabla\cdot\bm{J},
    \label{eq:app_conservation_law}
\end{equation}
with probability current
\begin{equation}
    J_i
    =
    F_i p
    -
    \partial_j(D_{ij}p).
    \label{eq:app_probability_current}
\end{equation}

Suppose now that $p_{\mathrm{ss}}$ is an invariant density for \eqref{eq:app_general_sde}. Then
\begin{equation}
    \mathcal{L}^*p_{\mathrm{ss}}=0,
    \qquad\text{or equivalently}\qquad
    \partial_i J_i^{\mathrm{ss}}=0,
    \label{eq:app_stationary_current_divfree}
\end{equation}
where
\begin{equation}
    J_i^{\mathrm{ss}}
    =
    F_i p_{\mathrm{ss}}
    -
    \partial_j(D_{ij}p_{\mathrm{ss}}).
    \label{eq:app_stationary_current}
\end{equation}
Thus the stationary probability current is divergence-free. Under the boundary and topological assumptions stated above, any such divergence-free current can be represented as the divergence of an antisymmetric tensor potential. More precisely, we assume that there exists a smooth matrix field $\bm{R}(\bm{x})$ satisfying
\begin{equation}
    R_{ij}(\bm{x})=-R_{ji}(\bm{x})
    \label{eq:app_R_antisymmetric}
\end{equation}
such that
\begin{equation}
    J_i^{\mathrm{ss}}(\bm{x})
    =
    \partial_j\!\left(
        R_{ij}(\bm{x})p_{\mathrm{ss}}(\bm{x})
    \right).
    \label{eq:app_current_R_representation}
\end{equation}
This is the antisymmetric-current representation of the stationary nonequilibrium flux. In Euclidean domains with sufficient decay it follows from the standard vector-potential representation of solenoidal vector fields. In domains with nontrivial topology, one must additionally exclude harmonic-current components or treat them separately.

Using \eqref{eq:app_stationary_current} and \eqref{eq:app_current_R_representation}, we obtain
\begin{equation}
    F_i p_{\mathrm{ss}}
    -
    \partial_j(D_{ij}p_{\mathrm{ss}})
    =
    \partial_j(R_{ij}p_{\mathrm{ss}}).
    \label{eq:app_current_balance}
\end{equation}
Therefore
\begin{equation}
    F_i p_{\mathrm{ss}}
    =
    \partial_j(D_{ij}p_{\mathrm{ss}})
    +
    \partial_j(R_{ij}p_{\mathrm{ss}})
    =
    \partial_j\!\left[
        (D_{ij}+R_{ij})p_{\mathrm{ss}}
    \right].
    \label{eq:app_Fp_as_divergence}
\end{equation}
We now define the mobility matrix
\begin{equation}
    M_{ij}(\bm{x})
    :=
    D_{ij}(\bm{x})+R_{ij}(\bm{x}),
    \label{eq:app_mobility_definition}
\end{equation}
or, in matrix form,
\begin{equation}
    \bm{M}(\bm{x})
    =
    \bm{D}(\bm{x})+\bm{R}(\bm{x}),
    \qquad
    \bm{D}(\bm{x})
    =
    \frac{\bm{M}(\bm{x})+\bm{M}(\bm{x})^T}{2},
    \qquad
    \bm{R}(\bm{x})
    =
    \frac{\bm{M}(\bm{x})-\bm{M}(\bm{x})^T}{2}.
    \label{eq:app_mobility_decomposition}
\end{equation}
Since $\bm{D}=\bm{\Sigma}\bm{\Sigma}^T$, the symmetric part of $\bm{M}$ is positive semidefinite.

Equation \eqref{eq:app_Fp_as_divergence} becomes
\begin{equation}
    F_i p_{\mathrm{ss}}
    =
    \partial_j(M_{ij}p_{\mathrm{ss}}).
    \label{eq:app_Fp_M}
\end{equation}
Expanding the derivative gives
\begin{equation}
    F_i p_{\mathrm{ss}}
    =
    (\partial_j M_{ij})p_{\mathrm{ss}}
    +
    M_{ij}\partial_j p_{\mathrm{ss}}.
    \label{eq:app_expand_derivative}
\end{equation}
Since $p_{\mathrm{ss}}>0$, we divide by $p_{\mathrm{ss}}$ and introduce the score
\begin{equation}
    s_j(\bm{x})
    :=
    \partial_j\log p_{\mathrm{ss}}(\bm{x})
    =
    \frac{\partial_j p_{\mathrm{ss}}(\bm{x})}{p_{\mathrm{ss}}(\bm{x})}.
    \label{eq:app_score_definition}
\end{equation}
This yields
\begin{equation}
    F_i(\bm{x})
    =
    M_{ij}(\bm{x})s_j(\bm{x})
    +
    \partial_j M_{ij}(\bm{x}).
    \label{eq:app_drift_component_form}
\end{equation}
In vector notation,
\begin{equation}
    \bm{F}(\bm{x})
    =
    \bm{M}(\bm{x})\bm{s}(\bm{x})
    +
    \nabla\cdot\bm{M}(\bm{x}),
    \label{eq:app_drift_vector_form}
\end{equation}
where
\begin{equation}
    (\nabla\cdot\bm{M})_i
    :=
    \partial_j M_{ij}.
    \label{eq:app_divergence_convention}
\end{equation}

Substituting \eqref{eq:app_drift_vector_form} into \eqref{eq:app_general_sde}, we obtain the score-based representation
\begin{equation}
    \d \bm{x}(t)
    =
    \left[
        \bm{M}(\bm{x}(t))\bm{s}(\bm{x}(t))
        +
        \nabla\cdot\bm{M}(\bm{x}(t))
    \right]\d t
    +
    \sqrt{2}\,\bm{\Sigma}(\bm{x}(t))\,\d \bm{W}_t,
    \label{eq:app_score_based_sde}
\end{equation}
with
\begin{equation}
    \bm{\Sigma}(\bm{x})\bm{\Sigma}(\bm{x})^T
    =
    \frac{\bm{M}(\bm{x})+\bm{M}(\bm{x})^T}{2}.
    \label{eq:app_sigma_constraint}
\end{equation}

Conversely, suppose that a diffusion is written in the form
\eqref{eq:app_score_based_sde}, with $\bm{s}=\nabla\log p_{\mathrm{ss}}$ and with
$\bm{M}=\bm{D}+\bm{R}$, where $\bm{D}=\bm{D}^T\succeq 0$ and
$\bm{R}=-\bm{R}^T$. The stationary current associated with $p_{\mathrm{ss}}$ is
\begin{align}
    J_i^{\mathrm{ss}}
    &=
    \left(
        M_{ij}s_j+\partial_jM_{ij}
    \right)p_{\mathrm{ss}}
    -
    \partial_j(D_{ij}p_{\mathrm{ss}})
    \nonumber\\
    &=
    \left(
        D_{ij}s_j+R_{ij}s_j+\partial_jD_{ij}+\partial_jR_{ij}
    \right)p_{\mathrm{ss}}
    -
    \partial_j(D_{ij}p_{\mathrm{ss}})
    \nonumber\\
    &=
    R_{ij}\partial_j p_{\mathrm{ss}}
    +
    (\partial_jR_{ij})p_{\mathrm{ss}}
    \nonumber\\
    &=
    \partial_j(R_{ij}p_{\mathrm{ss}}).
    \label{eq:app_converse_current}
\end{align}
Taking the divergence,
\begin{equation}
    \partial_iJ_i^{\mathrm{ss}}
    =
    \partial_i\partial_j(R_{ij}p_{\mathrm{ss}}).
    \label{eq:app_converse_divergence}
\end{equation}
Because $R_{ij}p_{\mathrm{ss}}$ is antisymmetric in $(i,j)$ while
$\partial_i\partial_j$ is symmetric in $(i,j)$, the right-hand side vanishes:
\begin{equation}
    \partial_i\partial_j(R_{ij}p_{\mathrm{ss}})
    =
    0.
    \label{eq:app_antisymmetry_cancellation}
\end{equation}
Therefore
\begin{equation}
    \mathcal{L}^*p_{\mathrm{ss}}
    =
    -\nabla\cdot\bm{J}^{\mathrm{ss}}
    =
    0,
    \label{eq:app_stationarity_verified}
\end{equation}
so $p_{\mathrm{ss}}$ is invariant for the diffusion.

The decomposition has a direct physical interpretation. The symmetric part $\bm{D}$ controls irreversible diffusion and dissipation, while the antisymmetric part $\bm{R}$ generates stationary probability currents that preserve the invariant density. If $\bm{R}=0$, the stationary current vanishes and the dynamics satisfy detailed balance. If $\bm{R}\neq 0$, the same invariant density can coexist with nonzero stationary currents, allowing the model to represent nonequilibrium dynamics.

\section{Conditional-score identity for lagged correlation derivatives}
\label{app:conditional_score_correlation_identity}

In this appendix we derive the identity
\begin{equation}
    \dot{\bm{C}}_{m,n}(t)
    =
    -
    \left\langle
        \phi_m(\bm{x}(t))\,
        \bm{s}_{t|0}(\bm{x}(t)\mid \bm{x}(0))^T\,
        \bm{M}(\bm{x}(0))\,
        \nabla\phi_n(\bm{x}(0))^T
    \right\rangle ,
    \label{eq:app_central_conditional_score_identity}
\end{equation}
used in Section~\ref{subsec:mobility_from_correlation_constraints}. We assume throughout that the process is stationary with invariant density $p_{\mathrm{ss}}$, that $p_{\mathrm{ss}}$ is strictly positive and smooth on the domain $\Omega\subseteq\mathbb{R}^D$, and that all observables and coefficients are sufficiently regular for the differentiations and integrations by parts below to be justified. Boundary terms are assumed to vanish, either because $\Omega=\mathbb{R}^D$ and the relevant quantities decay sufficiently fast, or because no-flux/periodic boundary conditions are imposed.

Consider the score-based diffusion
\begin{equation}
    \d \bm{x}(t)
    =
    \left[
        \bm{M}(\bm{x}(t))\bm{s}(\bm{x}(t))
        +
        \nabla\cdot\bm{M}(\bm{x}(t))
    \right]\d t
    +
    \sqrt{2}\,\bm{\Sigma}(\bm{x}(t))\,\d \bm{W}_t,
    \label{eq:app_score_based_sde_for_identity}
\end{equation}
where
\begin{equation}
    \bm{s}(\bm{x})
    =
    \nabla\log p_{\mathrm{ss}}(\bm{x}),
    \qquad
    \bm{\Sigma}(\bm{x})\bm{\Sigma}(\bm{x})^T
    =
    \bm{D}(\bm{x})
    =
    \frac{\bm{M}(\bm{x})+\bm{M}(\bm{x})^T}{2}.
    \label{eq:app_score_sde_definitions}
\end{equation}
The associated backward generator acts on smooth scalar functions $f$ as
\begin{equation}
    \mathcal{L}f
    =
    \left[
        M_{ij}s_j
        +
        \partial_j M_{ij}
    \right]\partial_i f
    +
    D_{ij}\partial_i\partial_j f ,
    \label{eq:app_backward_generator_score_form}
\end{equation}
where repeated indices are summed over. Equivalently, using
\begin{equation}
    M_{ij}
    =
    D_{ij}+R_{ij},
    \qquad
    R_{ij}=-R_{ji},
    \label{eq:app_M_D_R_decomposition_identity}
\end{equation}
and the stationarity of $p_{\mathrm{ss}}$, the generator admits the divergence-form adjoint relation
\begin{equation}
    \left\langle
        \mathcal{L}f\,h
    \right\rangle_{p_{\mathrm{ss}}}
    =
    -
    \left\langle
        \nabla f\,
        \bm{M}\,
        \nabla h^T
    \right\rangle_{p_{\mathrm{ss}}}
    +
    \text{terms antisymmetric under }f\leftrightarrow h,
    \label{eq:app_generator_weak_form_intro}
\end{equation}
which we now derive in the form needed for correlation functions.

Let
\begin{equation}
    \phi_m:\Omega\to\mathbb{R}^{d_m},
    \qquad
    \phi_n:\Omega\to\mathbb{R}^{d_n}
\end{equation}
be smooth vector-valued observables, and define
\begin{equation}
    \bm{C}_{m,n}(t)
    :=
    \mathbb{E}
    \left[
        \phi_m(\bm{x}(t))\,
        \phi_n(\bm{x}(0))^T
    \right].
    \label{eq:app_Cmn_definition}
\end{equation}
Let $K_t$ denote the Koopman semigroup of the diffusion,
\begin{equation}
    (K_t\phi_m)(\bm{x}_0)
    :=
    \mathbb{E}
    \left[
        \phi_m(\bm{x}(t))
        \,\middle|\,
        \bm{x}(0)=\bm{x}_0
    \right].
    \label{eq:app_Kt_definition}
\end{equation}
By stationarity,
\begin{equation}
    \bm{C}_{m,n}(t)
    =
    \int_\Omega
        (K_t\phi_m)(\bm{x}_0)\,
        \phi_n(\bm{x}_0)^T\,
        p_{\mathrm{ss}}(\bm{x}_0)\,
        \d\bm{x}_0 .
    \label{eq:app_Cmn_semigroup_form}
\end{equation}
Differentiating with respect to $t$ gives
\begin{equation}
    \dot{\bm{C}}_{m,n}(t)
    =
    \int_\Omega
        \mathcal{L}(K_t\phi_m)(\bm{x}_0)\,
        \phi_n(\bm{x}_0)^T\,
        p_{\mathrm{ss}}(\bm{x}_0)\,
        \d\bm{x}_0 ,
    \label{eq:app_Cmn_derivative_generator}
\end{equation}
where $\mathcal{L}$ acts componentwise on $K_t\phi_m$.

We now compute the $(a,b)$ entry of \eqref{eq:app_Cmn_derivative_generator}. Write
\[
    u(\bm{x})=(K_t\phi_m)_a(\bm{x}),
    \qquad
    v(\bm{x})=[\phi_n(\bm{x})]_b .
\]
Then
\begin{equation}
    [\dot{\bm{C}}_{m,n}(t)]_{ab}
    =
    \int_\Omega
        \mathcal{L}u(\bm{x})\,
        v(\bm{x})\,
        p_{\mathrm{ss}}(\bm{x})\,
        \d\bm{x}.
    \label{eq:app_C_entry_uv}
\end{equation}
Using the score-based drift, this becomes
\begin{equation}
\begin{aligned}
    \relax[\dot{\bm{C}}_{m,n}(t)]_{ab}
    &=
    \int_\Omega
    \left[
        \left(
            M_{ij}s_j+\partial_jM_{ij}
        \right)\partial_i u
        +
        D_{ij}\partial_i\partial_j u
    \right]
    v\,p_{\mathrm{ss}}\,\d\bm{x}.
\end{aligned}
\label{eq:app_generator_expanded_uv}
\end{equation}
Since $s_jp_{\mathrm{ss}}=\partial_jp_{\mathrm{ss}}$, the drift part can be written as
\begin{equation}
    \left(
        M_{ij}s_j+\partial_jM_{ij}
    \right)p_{\mathrm{ss}}
    =
    \partial_j(M_{ij}p_{\mathrm{ss}}).
    \label{eq:app_drift_divergence_identity}
\end{equation}
Therefore
\begin{equation}
\begin{aligned}
    \relax[\dot{\bm{C}}_{m,n}(t)]_{ab}
    &=
    \int_\Omega
        \partial_j(M_{ij}p_{\mathrm{ss}})
        \,\partial_i u\,v\,\d\bm{x}
    +
    \int_\Omega
        D_{ij}p_{\mathrm{ss}}
        \,\partial_i\partial_j u\,v\,\d\bm{x}.
\end{aligned}
\label{eq:app_two_terms_before_parts}
\end{equation}
Integrating the first term by parts in $x_j$ gives
\begin{equation}
    \int_\Omega
        \partial_j(M_{ij}p_{\mathrm{ss}})
        \,\partial_i u\,v\,\d\bm{x}
    =
    -
    \int_\Omega
        M_{ij}p_{\mathrm{ss}}
        \left[
            \partial_j\partial_i u\,v
            +
            \partial_i u\,\partial_j v
        \right]\d\bm{x}.
    \label{eq:app_first_term_parts}
\end{equation}
For the second term, use the symmetry of $\bm{D}$:
\begin{equation}
    \int_\Omega
        D_{ij}p_{\mathrm{ss}}
        \,\partial_i\partial_j u\,v\,\d\bm{x}
    =
    \int_\Omega
        D_{ij}p_{\mathrm{ss}}
        \,\partial_j\partial_i u\,v\,\d\bm{x}.
    \label{eq:app_second_term_symmetry}
\end{equation}
Substituting $\bm{M}=\bm{D}+\bm{R}$ into \eqref{eq:app_first_term_parts}, the terms involving $\partial_i\partial_j u$ cancel as follows:
\begin{equation}
    -
    M_{ij}\partial_j\partial_i u
    +
    D_{ij}\partial_j\partial_i u
    =
    -
    R_{ij}\partial_j\partial_i u
    =
    0,
    \label{eq:app_antisymmetric_hessian_cancellation}
\end{equation}
because $R_{ij}$ is antisymmetric while $\partial_j\partial_i u$ is symmetric. Hence
\begin{equation}
    [\dot{\bm{C}}_{m,n}(t)]_{ab}
    =
    -
    \int_\Omega
        \partial_i u(\bm{x})\,
        M_{ij}(\bm{x})\,
        \partial_j v(\bm{x})\,
        p_{\mathrm{ss}}(\bm{x})\,
        \d\bm{x}.
    \label{eq:app_C_entry_gradient_form}
\end{equation}
Returning to vector-valued notation, with
\[
    \nabla(K_t\phi_m)(\bm{x})\in\mathbb{R}^{d_m\times D},
    \qquad
    \nabla\phi_n(\bm{x})\in\mathbb{R}^{d_n\times D},
\]
we obtain
\begin{equation}
    \dot{\bm{C}}_{m,n}(t)
    =
    -
    \left\langle
        \nabla(K_t\phi_m)(\bm{x}_0)\,
        \bm{M}(\bm{x}_0)\,
        \nabla\phi_n(\bm{x}_0)^T
    \right\rangle_{p_{\mathrm{ss}}}.
    \label{eq:app_Cdot_gradient_Koopman_form}
\end{equation}

We next express $\nabla(K_t\phi_m)$ through the conditional score. Let
\begin{equation}
    p_t(\bm{x}_t\mid\bm{x}_0)
\end{equation}
be the transition density at lag $t$. Then
\begin{equation}
    (K_t\phi_m)(\bm{x}_0)
    =
    \int_\Omega
        \phi_m(\bm{x}_t)\,
        p_t(\bm{x}_t\mid\bm{x}_0)\,
        \d\bm{x}_t .
    \label{eq:app_Kt_transition_density}
\end{equation}
Assuming that differentiation under the integral is justified,
\begin{equation}
\begin{aligned}
    \nabla_{\bm{x}_0}(K_t\phi_m)(\bm{x}_0)
    &=
    \int_\Omega
        \phi_m(\bm{x}_t)\,
        \nabla_{\bm{x}_0}p_t(\bm{x}_t\mid\bm{x}_0)\,
        \d\bm{x}_t \\
    &=
    \int_\Omega
        \phi_m(\bm{x}_t)\,
        p_t(\bm{x}_t\mid\bm{x}_0)\,
        \nabla_{\bm{x}_0}
        \log p_t(\bm{x}_t\mid\bm{x}_0)^T
        \d\bm{x}_t .
\end{aligned}
\label{eq:app_grad_Kt_transition_score}
\end{equation}
Define the conditional transition score by
\begin{equation}
    \bm{s}_{t|0}(\bm{x}_t\mid\bm{x}_0)
    :=
    \nabla_{\bm{x}_0}
    \log p_t(\bm{x}_t\mid\bm{x}_0)
    \in\mathbb{R}^{D}.
    \label{eq:app_conditional_score_definition}
\end{equation}
Then
\begin{equation}
    \nabla_{\bm{x}_0}(K_t\phi_m)(\bm{x}_0)
    =
    \mathbb{E}
    \left[
        \phi_m(\bm{x}_t)\,
        \bm{s}_{t|0}(\bm{x}_t\mid\bm{x}_0)^T
        \,\middle|\,
        \bm{x}_0
    \right].
    \label{eq:app_grad_Kt_conditional_expectation}
\end{equation}
Substituting \eqref{eq:app_grad_Kt_conditional_expectation} into
\eqref{eq:app_Cdot_gradient_Koopman_form} gives
\begin{equation}
\begin{aligned}
    \dot{\bm{C}}_{m,n}(t)
    &=
    -
    \int_\Omega
    \left[
        \int_\Omega
            \phi_m(\bm{x}_t)\,
            \bm{s}_{t|0}(\bm{x}_t\mid\bm{x}_0)^T\,
            p_t(\bm{x}_t\mid\bm{x}_0)
            \d\bm{x}_t
    \right]
    \bm{M}(\bm{x}_0)
    \nabla\phi_n(\bm{x}_0)^T
    p_{\mathrm{ss}}(\bm{x}_0)
    \d\bm{x}_0 .
\end{aligned}
\label{eq:app_substitute_conditional_score}
\end{equation}
Recognizing the joint stationary law
\begin{equation}
    p_{\mathrm{ss}}(\bm{x}_0)
    p_t(\bm{x}_t\mid\bm{x}_0)
\end{equation}
of the pair $(\bm{x}(0),\bm{x}(t))$, we obtain
\begin{equation}
    \dot{\bm{C}}_{m,n}(t)
    =
    -
    \left\langle
        \phi_m(\bm{x}(t))\,
        \bm{s}_{t|0}(\bm{x}(t)\mid \bm{x}(0))^T\,
        \bm{M}(\bm{x}(0))\,
        \nabla\phi_n(\bm{x}(0))^T
    \right\rangle ,
    \label{eq:app_final_conditional_score_identity}
\end{equation}
which is \eqref{eq:app_central_conditional_score_identity}.

This identity has three important consequences. First, it expresses the action of the generator on finite-lag correlations without estimating instantaneous time derivatives from trajectory increments. Second, it depends on the transition density only through its score with respect to the initial condition, which can be estimated from lagged pairs by score-matching methods. Third, the identity is linear in the unknown mobility field $\bm{M}$, making it suitable for inverse problems in which $\bm{M}$ is represented by a finite-dimensional basis or by a neural network.

\section{Score estimation by denoising score matching}
\label{app:score_estimation_dsm}

This appendix describes how the stationary score and the conditional transition score used in the main text can be estimated from data by denoising score matching. The first subsection treats the stationary score
\[
    \bm{s}(\bm{x})
    =
    \nabla_{\bm{x}}\log p_{\mathrm{ss}}(\bm{x}),
\]
while the second subsection uses the same construction to estimate the conditional score
\[
    \bm{s}_{t|0}(\bm{x}_t\mid\bm{x}_0)
    =
    \nabla_{\bm{x}_0}\log p_t(\bm{x}_t\mid\bm{x}_0).
\]

\subsection{Estimating the stationary score}
\label{app:stationary_score_dsm}

Let $\bm{x}\sim p_{\mathrm{ss}}$ denote a sample from the stationary distribution of the resolved process. The score
\begin{equation}
    \bm{s}(\bm{x})
    :=
    \nabla_{\bm{x}}\log p_{\mathrm{ss}}(\bm{x})
    \label{eq:app_stationary_score_def}
\end{equation}
is the central object needed to impose the observed invariant density in the score-based Langevin representation. Directly estimating \eqref{eq:app_stationary_score_def} from data is difficult because the density $p_{\mathrm{ss}}$ is unknown, especially in high dimension. Denoising score matching avoids explicit density estimation by learning the score of a noise-regularized version of the data distribution.

For a noise level $\sigma>0$, define the Gaussian-corrupted variable
\begin{equation}
    \widetilde{\bm{x}}
    =
    \bm{x}
    +
    \sigma\bm{\xi},
    \qquad
    \bm{\xi}\sim\mathcal{N}(\bm{0},\bm{I}_D).
    \label{eq:app_gaussian_corruption_stationary}
\end{equation}
The density of $\widetilde{\bm{x}}$ is the Gaussian convolution
\begin{equation}
    p_\sigma(\widetilde{\bm{x}})
    =
    \int_{\mathbb{R}^D}
        p_{\mathrm{ss}}(\bm{x})\,
        \varphi_\sigma(\widetilde{\bm{x}}-\bm{x})\,
        \d\bm{x},
    \label{eq:app_smoothed_density}
\end{equation}
where
\begin{equation}
    \varphi_\sigma(\bm{z})
    =
    (2\pi\sigma^2)^{-D/2}
    \exp\left(
        -\frac{\|\bm{z}\|^2}{2\sigma^2}
    \right).
    \label{eq:app_gaussian_kernel}
\end{equation}
The corresponding smoothed score is
\begin{equation}
    \bm{s}_\sigma(\widetilde{\bm{x}})
    :=
    \nabla_{\widetilde{\bm{x}}}
    \log p_\sigma(\widetilde{\bm{x}}).
    \label{eq:app_smoothed_score}
\end{equation}

The key identity behind denoising score matching is that this score can be written as a conditional expectation involving the clean sample $\bm{x}$. Indeed,
\begin{align}
    \nabla_{\widetilde{\bm{x}}}p_\sigma(\widetilde{\bm{x}})
    &=
    \int
        p_{\mathrm{ss}}(\bm{x})\,
        \nabla_{\widetilde{\bm{x}}}
        \varphi_\sigma(\widetilde{\bm{x}}-\bm{x})
        \,\d\bm{x}
    \nonumber\\
    &=
    \int
        p_{\mathrm{ss}}(\bm{x})\,
        \left(
            -\frac{\widetilde{\bm{x}}-\bm{x}}{\sigma^2}
        \right)
        \varphi_\sigma(\widetilde{\bm{x}}-\bm{x})
        \,\d\bm{x}.
    \label{eq:app_grad_smoothed_density}
\end{align}
Dividing by $p_\sigma(\widetilde{\bm{x}})$ gives
\begin{equation}
    \bm{s}_\sigma(\widetilde{\bm{x}})
    =
    \mathbb{E}
    \left[
        -\frac{\widetilde{\bm{x}}-\bm{x}}{\sigma^2}
        \,\middle|\,
        \widetilde{\bm{x}}
    \right].
    \label{eq:app_denoising_identity_stationary}
\end{equation}
Thus the score of the smoothed density is the optimal least-squares predictor of the denoising vector
\begin{equation}
    -\frac{\widetilde{\bm{x}}-\bm{x}}{\sigma^2}.
\end{equation}

Let $\bm{s}_\theta(\widetilde{\bm{x}},\sigma)$ be a parametric model for the smoothed score. For a fixed noise level $\sigma$, denoising score matching minimizes
\begin{equation}
    \mathcal{L}_{\mathrm{DSM}}(\theta;\sigma)
    =
    \mathbb{E}_{\bm{x}\sim p_{\mathrm{ss}}}
    \mathbb{E}_{\bm{\xi}\sim\mathcal{N}(\bm{0},\bm{I}_D)}
    \left[
        \left\|
            \bm{s}_\theta(\bm{x}+\sigma\bm{\xi},\sigma)
            +
            \frac{\sigma\bm{\xi}}{\sigma^2}
        \right\|^2
    \right].
    \label{eq:app_dsm_loss_fixed_sigma}
\end{equation}
Equivalently, using $\widetilde{\bm{x}}=\bm{x}+\sigma\bm{\xi}$,
\begin{equation}
    \mathcal{L}_{\mathrm{DSM}}(\theta;\sigma)
    =
    \mathbb{E}
    \left[
        \left\|
            \bm{s}_\theta(\widetilde{\bm{x}},\sigma)
            +
            \frac{\widetilde{\bm{x}}-\bm{x}}{\sigma^2}
        \right\|^2
    \right].
    \label{eq:app_dsm_loss_fixed_sigma_tilde}
\end{equation}
The minimizer satisfies
\begin{equation}
    \bm{s}_{\theta^\ast}(\widetilde{\bm{x}},\sigma)
    =
    \nabla_{\widetilde{\bm{x}}}
    \log p_\sigma(\widetilde{\bm{x}})
    \qquad
    p_\sigma\text{-a.e.}
    \label{eq:app_dsm_optimality}
\end{equation}
This follows directly from the projection property of conditional expectation: the least-squares optimal predictor of
$-(\widetilde{\bm{x}}-\bm{x})/\sigma^2$ given $\widetilde{\bm{x}}$ is precisely the conditional expectation in \eqref{eq:app_denoising_identity_stationary}.

In practice, one usually trains over a distribution of noise levels. Let $\sigma\sim\rho(\sigma)$, with support contained in $[\sigma_{\min},\sigma_{\max}]$. The multi-noise denoising score matching objective is
\begin{equation}
    \mathcal{L}_{\mathrm{DSM}}(\theta)
    =
    \mathbb{E}_{\sigma\sim\rho}
    \mathbb{E}_{\bm{x}\sim p_{\mathrm{ss}}}
    \mathbb{E}_{\bm{\xi}\sim\mathcal{N}(\bm{0},\bm{I}_D)}
    \left[
        \omega(\sigma)
        \left\|
            \bm{s}_\theta(\bm{x}+\sigma\bm{\xi},\sigma)
            +
            \frac{\bm{\xi}}{\sigma}
        \right\|^2
    \right],
    \label{eq:app_multinoise_dsm_loss}
\end{equation}
where $\omega(\sigma)>0$ is a weighting function. For each fixed $\sigma$, the population minimizer is the smoothed score $\nabla\log p_\sigma$. The desired stationary score is then obtained in the small-noise limit:
\begin{equation}
    \bm{s}(\bm{x})
    =
    \nabla\log p_{\mathrm{ss}}(\bm{x})
    =
    \lim_{\sigma\to 0}
    \nabla\log p_\sigma(\bm{x}).
    \label{eq:app_score_small_noise_limit}
\end{equation}
With finite data one uses a small but nonzero $\sigma_{\min}$, which regularizes the score estimate and avoids fitting sampling noise.

Given stationary samples $\{\bm{x}^{(k)}\}_{k=1}^N$, the empirical version of \eqref{eq:app_multinoise_dsm_loss} is
\begin{equation}
    \widehat{\mathcal{L}}_{\mathrm{DSM}}(\theta)
    =
    \frac{1}{N}
    \sum_{k=1}^N
    \mathbb{E}_{\sigma,\bm{\xi}}
    \left[
        \omega(\sigma)
        \left\|
            \bm{s}_\theta(\bm{x}^{(k)}+\sigma\bm{\xi},\sigma)
            +
            \frac{\bm{\xi}}{\sigma}
        \right\|^2
    \right].
    \label{eq:app_empirical_stationary_dsm}
\end{equation}
After training, the stationary score used in the reduced model is approximated by
\begin{equation}
    \widehat{\bm{s}}(\bm{x})
    \approx
    \bm{s}_\theta(\bm{x},\sigma_{\min}),
    \label{eq:app_stationary_score_estimator}
\end{equation}
or by an average over a small range of low noise levels.

\subsection{Estimating the conditional transition score from joint scores}
\label{app:conditional_score_dsm}

We now describe how the conditional transition score appearing in the
correlation identity is estimated in the manuscript. The required object is
\begin{equation}
    \bm{s}_{t|0}(\bm{x}_t\mid\bm{x}_0)
    =
    \nabla_{\bm{x}_0}
    \log p_t(\bm{x}_t\mid\bm{x}_0).
    \label{eq:app_transition_score_def}
\end{equation}
This score is a derivative with respect to the initial state, not the final
state. It measures the sensitivity of the finite-lag transition density to
perturbations of the initial condition $\bm{x}_0$.

The data provide stationary lagged pairs
\begin{equation}
    (\bm{x}_0,\bm{x}_t)
    \sim
    p_{0,t}(\bm{x}_0,\bm{x}_t)
    =
    p_{\mathrm{ss}}(\bm{x}_0)\,
    p_t(\bm{x}_t\mid\bm{x}_0),
    \qquad
    t\in\mathcal{T}.
    \label{eq:app_lagged_pair_distribution}
\end{equation}
For each lag $t$, we regard the observed pairs $(\bm{x}_0,\bm{x}_t)$ as
samples from the joint density $p_{0,t}$. We then estimate the joint score
\begin{equation}
    \bm{s}_{0,t}(\bm{x}_0,\bm{x}_t)
    :=
    \nabla_{(\bm{x}_0,\bm{x}_t)}
    \log p_{0,t}(\bm{x}_0,\bm{x}_t).
    \label{eq:app_joint_score_def}
\end{equation}
Writing this score in block form gives
\begin{equation}
    \bm{s}_{0,t}(\bm{x}_0,\bm{x}_t)
    =
    \begin{pmatrix}
        \bm{s}_{0,t}^{(0)}(\bm{x}_0,\bm{x}_t) \\
        \bm{s}_{0,t}^{(t)}(\bm{x}_0,\bm{x}_t)
    \end{pmatrix},
    \label{eq:app_joint_score_blocks}
\end{equation}
where
\begin{equation}
    \bm{s}_{0,t}^{(0)}(\bm{x}_0,\bm{x}_t)
    =
    \nabla_{\bm{x}_0}
    \log p_{0,t}(\bm{x}_0,\bm{x}_t),
    \qquad
    \bm{s}_{0,t}^{(t)}(\bm{x}_0,\bm{x}_t)
    =
    \nabla_{\bm{x}_t}
    \log p_{0,t}(\bm{x}_0,\bm{x}_t).
    \label{eq:app_joint_score_components}
\end{equation}

The conditional transition density satisfies
\begin{equation}
    p_t(\bm{x}_t\mid\bm{x}_0)
    =
    \frac{
        p_{0,t}(\bm{x}_0,\bm{x}_t)
    }{
        p_{\mathrm{ss}}(\bm{x}_0)
    }.
    \label{eq:app_transition_density_from_joint}
\end{equation}
Taking the logarithm and differentiating with respect to $\bm{x}_0$ gives
\begin{equation}
    \nabla_{\bm{x}_0}\log p_t(\bm{x}_t\mid\bm{x}_0)
    =
    \nabla_{\bm{x}_0}
    \log p_{0,t}(\bm{x}_0,\bm{x}_t)
    -
    \nabla_{\bm{x}_0}
    \log p_{\mathrm{ss}}(\bm{x}_0).
    \label{eq:app_transition_score_from_joint}
\end{equation}
Therefore,
\begin{equation}
    \bm{s}_{t|0}(\bm{x}_t\mid\bm{x}_0)
    =
    \bm{s}_{0,t}^{(0)}(\bm{x}_0,\bm{x}_t)
    -
    \bm{s}(\bm{x}_0),
    \label{eq:app_conditional_score_joint_minus_stationary}
\end{equation}
where
\begin{equation}
    \bm{s}(\bm{x}_0)
    :=
    \nabla_{\bm{x}_0}\log p_{\mathrm{ss}}(\bm{x}_0)
    \label{eq:app_stationary_score_in_conditional_section}
\end{equation}
is the stationary score. Thus, the conditional transition score is obtained
from the $\bm{x}_0$-component of the joint score by subtracting the stationary
score.

In practice, the joint score is learned by denoising score matching on the
augmented variable
\begin{equation}
    \bm{z}_t
    =
    \begin{pmatrix}
        \bm{x}_0 \\
        \bm{x}_t
    \end{pmatrix}
    \in \mathbb{R}^{2D}.
    \label{eq:app_augmented_lagged_variable}
\end{equation}
For a fixed lag $t$, define the corrupted pair
\begin{equation}
    \widetilde{\bm{z}}_t
    =
    \bm{z}_t
    +
    \sigma \bm{\xi},
    \qquad
    \bm{\xi}\sim\mathcal{N}(\bm{0},\bm{I}_{2D}).
    \label{eq:app_joint_corruption}
\end{equation}
Equivalently,
\begin{equation}
    \widetilde{\bm{x}}_0
    =
    \bm{x}_0+\sigma\bm{\xi}_0,
    \qquad
    \widetilde{\bm{x}}_t
    =
    \bm{x}_t+\sigma\bm{\xi}_t,
    \qquad
    \bm{\xi}_0,\bm{\xi}_t\sim\mathcal{N}(\bm{0},\bm{I}_D).
    \label{eq:app_joint_corruption_blocks}
\end{equation}
Let $p_{\sigma,0,t}$ denote the Gaussian-smoothed joint density of
$(\bm{x}_0,\bm{x}_t)$. The corresponding smoothed joint score is
\begin{equation}
    \bm{s}_{\sigma,0,t}(\widetilde{\bm{x}}_0,\widetilde{\bm{x}}_t)
    :=
    \nabla_{(\widetilde{\bm{x}}_0,\widetilde{\bm{x}}_t)}
    \log p_{\sigma,0,t}(\widetilde{\bm{x}}_0,\widetilde{\bm{x}}_t).
    \label{eq:app_smoothed_joint_score}
\end{equation}
By the standard denoising identity,
\begin{equation}
    \bm{s}_{\sigma,0,t}(\widetilde{\bm{z}}_t)
    =
    \mathbb{E}
    \left[
        -\frac{\widetilde{\bm{z}}_t-\bm{z}_t}{\sigma^2}
        \,\middle|\,
        \widetilde{\bm{z}}_t
    \right].
    \label{eq:app_joint_denoising_identity}
\end{equation}

A neural score model
\[
    \bm{q}_{\psi}(\widetilde{\bm{x}}_0,\widetilde{\bm{x}}_t,t,\sigma)
    \in \mathbb{R}^{2D}
\]
is trained to approximate this smoothed joint score by minimizing
\begin{equation}
    \mathcal{L}_{\mathrm{jDSM}}(\psi)
    =
    \mathbb{E}_{t\sim\rho_{\mathcal{T}}}
    \mathbb{E}_{(\bm{x}_0,\bm{x}_t)}
    \mathbb{E}_{\sigma\sim\rho}
    \mathbb{E}_{\bm{\xi}\sim\mathcal{N}(\bm{0},\bm{I}_{2D})}
    \left[
        \omega(\sigma)
        \left\|
            \bm{q}_{\psi}
            \left(
                \bm{x}_0+\sigma\bm{\xi}_0,
                \bm{x}_t+\sigma\bm{\xi}_t,
                t,
                \sigma
            \right)
            +
            \frac{\bm{\xi}}{\sigma}
        \right\|^2
    \right],
    \label{eq:app_joint_dsm_loss}
\end{equation}
where $\bm{\xi}=(\bm{\xi}_0,\bm{\xi}_t)$ and
$\rho_{\mathcal{T}}$ is the sampling distribution over the selected lag times.
For each fixed $(t,\sigma)$, the population minimizer satisfies
\begin{equation}
    \bm{q}_{\psi^\ast}
    (\widetilde{\bm{x}}_0,\widetilde{\bm{x}}_t,t,\sigma)
    =
    \nabla_{(\widetilde{\bm{x}}_0,\widetilde{\bm{x}}_t)}
    \log p_{\sigma,0,t}
    (\widetilde{\bm{x}}_0,\widetilde{\bm{x}}_t).
    \label{eq:app_joint_dsm_optimality}
\end{equation}

We write the learned joint score in block form as
\begin{equation}
    \bm{q}_{\psi}
    (\widetilde{\bm{x}}_0,\widetilde{\bm{x}}_t,t,\sigma)
    =
    \begin{pmatrix}
        \bm{q}_{\psi}^{(0)}
        (\widetilde{\bm{x}}_0,\widetilde{\bm{x}}_t,t,\sigma)
        \\
        \bm{q}_{\psi}^{(t)}
        (\widetilde{\bm{x}}_0,\widetilde{\bm{x}}_t,t,\sigma)
    \end{pmatrix}.
    \label{eq:app_learned_joint_score_blocks}
\end{equation}
The component needed for the correlation identity is the derivative with
respect to the initial state. Combining the learned joint score with the
stationary score estimator from Section~\ref{app:stationary_score_dsm}, we
therefore use
\begin{equation}
    \widehat{\bm{s}}_{t|0}(\bm{x}_t\mid\bm{x}_0)
    =
    \bm{q}_{\psi}^{(0)}
    (\bm{x}_0,\bm{x}_t,t,\sigma_{\min})
    -
    \bm{s}_{\theta}(\bm{x}_0,\sigma_{\min}).
    \label{eq:app_conditional_score_estimator}
\end{equation}
In the limit of infinite data, infinite model capacity, and
$\sigma_{\min}\to0$, this estimator converges to
\begin{equation}
    \widehat{\bm{s}}_{t|0}(\bm{x}_t\mid\bm{x}_0)
    \longrightarrow
    \nabla_{\bm{x}_0}
    \log p_t(\bm{x}_t\mid\bm{x}_0).
    \label{eq:app_conditional_score_convergence}
\end{equation}

Given a trajectory sampled at times $\{t_k\}$, empirical training pairs are
constructed by selecting
\begin{equation}
    (\bm{x}_0^{(k)},\bm{x}_t^{(k)})
    =
    (\bm{x}(t_k),\bm{x}(t_k+t))
\end{equation}
for each lag $t\in\mathcal{T}$. If the data are unevenly sampled, one may
instead use all pairs whose time separation lies in a prescribed lag window.
The empirical joint denoising objective corresponding to
\eqref{eq:app_joint_dsm_loss} is
\begin{equation}
    \widehat{\mathcal{L}}_{\mathrm{jDSM}}(\psi)
    =
    \frac{1}{|\mathcal{P}|}
    \sum_{(k,t)\in\mathcal{P}}
    \mathbb{E}_{\sigma,\bm{\xi}}
    \left[
        \omega(\sigma)
        \left\|
            \bm{q}_{\psi}
            \left(
                \bm{x}_0^{(k)}+\sigma\bm{\xi}_0,
                \bm{x}_t^{(k)}+\sigma\bm{\xi}_t,
                t,
                \sigma
            \right)
            +
            \frac{\bm{\xi}}{\sigma}
        \right\|^2
    \right],
    \label{eq:app_empirical_joint_dsm}
\end{equation}
where $\mathcal{P}$ denotes the collection of observed lagged pairs used for
training. The lag $t$ can be supplied to the network as an additional input,
allowing a single model to learn joint scores over multiple time separations.


\section{Mean mobility decomposition and constant-closure limit}
\label{app:mean_mobility_correction_derivation}

This appendix derives the identities used in
Section~\ref{subsec:mean_mobility_and_corrections}. We start from the
correlation identity established in
Appendix~\ref{app:conditional_score_correlation_identity} and then specialize it
to the mean--fluctuation decomposition of the mobility. The notation follows
Sec.~\ref{subsec:gfdt_interpretation_correlation_constraint}: for a perturbing
drift \(\bm{b}\), the conjugate variable is \(B_{\bm{b}}\). When an observable
\(\phi_n\) is involved, \(\nabla\phi_{n,j}\) enters through the componentwise
drift field \(\bm{b}_{\bm{A},n,j}=-\bm{A}\nabla\phi_{n,j}\), not through a new
conjugate-variable definition.

Let $\phi_m:\mathbb{R}^D\to\mathbb{R}^{d_m}$ and $\phi_n:\mathbb{R}^D\to\mathbb{R}^{d_n}$ be smooth observables, and define
\begin{equation}
    \bm{C}_{m,n}(t)
    =
    \left\langle
        \phi_m(\bm{x}(t))\,
        \phi_n(\bm{x}(0))^T
    \right\rangle .
    \label{eq:app_mean_Cmn_def}
\end{equation}
From Appendix~\ref{app:conditional_score_correlation_identity}, the correlation derivative can be written equivalently as
\begin{equation}
    \dot{\bm{C}}_{m,n}(t)
    =
    -
    \left\langle
        \nabla(K_t\phi_m)(\bm{x}_0)\,
        \bm{M}(\bm{x}_0)\,
        \nabla\phi_n(\bm{x}_0)^T
    \right\rangle_{p_{\mathrm{ss}}},
    \label{eq:app_Cdot_gradient_form_start}
\end{equation}
where $K_t$ is the Koopman semigroup,
\begin{equation}
    (K_t\phi_m)(\bm{x}_0)
    =
    \mathbb{E}
    \left[
        \phi_m(\bm{x}(t))
        \,\middle|\,
        \bm{x}(0)=\bm{x}_0
    \right].
    \label{eq:app_Kt_phi_def_mean}
\end{equation}
Componentwise, if $u=(K_t\phi_m)_a$ and $v=(\phi_n)_b$, then
\begin{equation}
    [\dot{\bm{C}}_{m,n}(t)]_{ab}
    =
    -
    \int
        \partial_i u(\bm{x})\,
        M_{ij}(\bm{x})\,
        \partial_j v(\bm{x})\,
        p_{\mathrm{ss}}(\bm{x})\,\d\bm{x}.
    \label{eq:app_Cdot_component_start}
\end{equation}

We decompose the mobility as
\begin{equation}
    \bm{M}(\bm{x})
    =
    \bm{\Phi}
    +
    \delta\bm{M}(\bm{x}),
    \qquad
    \bm{\Phi}
    =
    \left\langle
        \bm{M}
    \right\rangle_{p_{\mathrm{ss}}},
    \qquad
    \left\langle
        \delta\bm{M}
    \right\rangle_{p_{\mathrm{ss}}}
    =
    \bm{0}.
    \label{eq:app_M_Phi_delta_def}
\end{equation}
Substituting \eqref{eq:app_M_Phi_delta_def} into \eqref{eq:app_Cdot_gradient_form_start} gives
\begin{equation}
\begin{aligned}
    \dot{\bm{C}}_{m,n}(t)
    &=
    -
    \left\langle
        \nabla(K_t\phi_m)(\bm{x}_0)\,
        \bm{\Phi}\,
        \nabla\phi_n(\bm{x}_0)^T
    \right\rangle
    \\
    &\hspace{0.5cm}
    -
    \left\langle
        \nabla(K_t\phi_m)(\bm{x}_0)\,
        \delta\bm{M}(\bm{x}_0)\,
        \nabla\phi_n(\bm{x}_0)^T
    \right\rangle .
\end{aligned}
\label{eq:app_Cdot_Phi_delta_gradient}
\end{equation}

We now rewrite each term as a GFDT correlation with the conjugate variable
defined in \eqref{eq:gfdt_conjugate_variable_vector}. For a matrix field
\(\bm{A}(\bm{x})\), use the componentwise drift field
\begin{equation}
    \bm{b}_{\bm{A},n,b}(\bm{x})
    =
    -
    \bm{A}(\bm{x})\nabla\phi_{n,b}(\bm{x}) .
    \label{eq:app_b_A_n_b_def}
\end{equation}
Then
\begin{equation}
    B_{\bm{b}_{\bm{A},n,b}}(\bm{x})
    =
    \frac{1}{p_{\mathrm{ss}}(\bm{x})}
    \partial_i
    \left[
        p_{\mathrm{ss}}(\bm{x})\,
        A_{ij}(\bm{x})\,
        \partial_j \phi_{n,b}(\bm{x})
    \right].
    \label{eq:app_B_b_A_component}
\end{equation}
Using integration by parts and assuming the boundary terms vanish,
\begin{align}
    &-
    \int
        \partial_i u(\bm{x})\,
        A_{ij}(\bm{x})\,
        \partial_j v(\bm{x})\,
        p_{\mathrm{ss}}(\bm{x})\,\d\bm{x}
    \nonumber\\
    &\hspace{1cm}
    =
    \int
        u(\bm{x})\,
        \partial_i
        \left[
            p_{\mathrm{ss}}(\bm{x})\,
            A_{ij}(\bm{x})\,
            \partial_j v(\bm{x})
        \right]
        \d\bm{x}
    \nonumber\\
    &\hspace{1cm}
    =
    \int
        u(\bm{x})\,
        B_{\bm{b}_{\bm{A},n,b}}(\bm{x})\,
        p_{\mathrm{ss}}(\bm{x})\,\d\bm{x}.
    \label{eq:app_integration_by_parts_gamma}
\end{align}
Since $u(\bm{x}_0)=(K_t\phi_m)_a(\bm{x}_0)$, the last expression can be written as a stationary lagged expectation:
\begin{equation}
    \int
        (K_t\phi_m)_a(\bm{x}_0)\,
        B_{\bm{b}_{\bm{A},n,b}}(\bm{x}_0)\,
        p_{\mathrm{ss}}(\bm{x}_0)\,\d\bm{x}_0
    =
    \left\langle
        [\phi_m(\bm{x}(t))]_a\,
        B_{\bm{b}_{\bm{A},n,b}}(\bm{x}(0))
    \right\rangle .
    \label{eq:app_semigroup_to_lagged_gamma}
\end{equation}
Applying this identity first with $\bm{A}=\bm{\Phi}$ and then with
$\bm{A}=\delta\bm{M}$ gives the corresponding decomposition into the constant
and correction impulse-response channels,
\begin{equation}
    \left[\dot{\bm{C}}_{m,n}(t)\right]_{:,b}
    =
    \left\langle
        \phi_m(\bm{x}(t))\,
        B_{\bm{b}_{\bm{\Phi},n,b}}(\bm{x}(0))
    \right\rangle
    +
    \left\langle
        \phi_m(\bm{x}(t))\,
        B_{\bm{b}_{\delta\bm{M},n,b}}(\bm{x}(0))
    \right\rangle,
    \qquad
    b=1,\dots,d_n .
    \label{eq:app_Cdot_gamma_Phi_delta_final}
\end{equation}
This is the decomposition used in \eqref{eq:Cdot_mean_delta_gamma_method}.

We next show that $\bm{\Phi}$ is determined by the right derivative at the origin of the coordinate correlation. Let
\begin{equation}
    \bm{\iota}(\bm{x})=\bm{x}
    \label{eq:app_coordinate_observable_iota}
\end{equation}
be the identity observable. Then
\begin{equation}
    \nabla\bm{\iota}(\bm{x})=\bm{I}.
    \label{eq:app_coordinate_gradient_identity}
\end{equation}
Taking $\phi_m=\phi_n=\bm{\iota}$ in \eqref{eq:app_Cdot_gradient_form_start} yields
\begin{equation}
    \dot{\bm{C}}_{\bm{x},\bm{x}}(t)
    =
    -
    \left\langle
        \nabla(K_t\bm{\iota})(\bm{x}_0)\,
        \bm{M}(\bm{x}_0)
    \right\rangle .
    \label{eq:app_coordinate_Cdot_general}
\end{equation}
At $t=0$, $K_0\bm{\iota}=\bm{\iota}$ and therefore $\nabla(K_0\bm{\iota})=\bm{I}$. Hence
\begin{equation}
    \dot{\bm{C}}_{\bm{x},\bm{x}}(0^+)
    =
    -
    \left\langle
        \bm{M}
    \right\rangle
    =
    -
    \bm{\Phi}.
    \label{eq:app_Cdot_zero_Phi}
\end{equation}
Equivalently,
\begin{equation}
    \bm{\Phi}
    =
    -
    \dot{\bm{C}}_{\bm{x},\bm{x}}(0^+).
    \label{eq:app_Phi_from_Cdot_zero}
\end{equation}
The right derivative is used because diffusion sample paths are not differentiable, while the semigroup correlation function has a well-defined right derivative under the generator assumptions used above.

Finally, we analyze the contribution of $\delta\bm{M}$ to coordinate correlations. Define the finite-lag conditional mean
\begin{equation}
    \bm{m}_t(\bm{x}_0)
    :=
    (K_t\bm{\iota})(\bm{x}_0)
    =
    \mathbb{E}
    \left[
        \bm{x}(t)
        \,\middle|\,
        \bm{x}(0)=\bm{x}_0
    \right].
    \label{eq:app_conditional_mean_def}
\end{equation}
The average Jacobian of this conditional mean can be written either with the
conditional transition score or with the stationary score. Indeed,
\begin{equation}
    \nabla\bm{m}_t(\bm{x}_0)
    =
    \mathbb{E}
    \left[
        \bm{x}_t\,
        \bm{s}_{t|0}(\bm{x}_t\mid\bm{x}_0)^T
        \,\middle|\,
        \bm{x}_0
    \right],
    \label{eq:app_conditional_mean_gradient_score}
\end{equation}
and therefore, using integration by parts against \(p_{\mathrm{ss}}\),
\begin{align}
    \left\langle
        \bm{x}_t\,
        \bm{s}_{t|0}(\bm{x}_t\mid\bm{x}_0)^T
    \right\rangle_{\mathrm{obs}}
    &=
    \left\langle
        \nabla\bm{m}_t(\bm{x}_0)
    \right\rangle_{p_{\mathrm{ss}}}
    \nonumber\\
    &=
    \int
        \nabla\bm{m}_t(\bm{x}_0)\,
        p_{\mathrm{ss}}(\bm{x}_0)\,
        \d\bm{x}_0
    \nonumber\\
    &=
    -
    \int
        \bm{m}_t(\bm{x}_0)\,
        \nabla p_{\mathrm{ss}}(\bm{x}_0)^T\,
        \d\bm{x}_0
    \nonumber\\
    &=
    -
    \left\langle
        \bm{x}_t\,
        \bm{s}(\bm{x}_0)^T
    \right\rangle_{\mathrm{obs}} .
    \label{eq:app_conditional_score_stationary_score_equivalence}
\end{align}
Here \(\bm{s}=\nabla\log p_{\mathrm{ss}}\) is the stationary score.
Then \eqref{eq:app_coordinate_Cdot_general}, together with the decomposition
$\bm{M}=\bm{\Phi}+\delta\bm{M}$ and
\eqref{eq:app_conditional_score_stationary_score_equivalence}, gives
\begin{align}
    \dot{\bm{C}}_{\bm{x},\bm{x}}(t)
    &=
    \left\langle
        \bm{x}_t\,
        \bm{s}(\bm{x}_0)^T
    \right\rangle_{\mathrm{obs}}
    \bm{\Phi}
    \\
    &\quad-
    \left\langle
        \nabla\bm{m}_t(\bm{x}_0)\,
        \delta\bm{M}(\bm{x}_0)
    \right\rangle .
    \label{eq:app_coordinate_Cdot_conditional_mean}
\end{align}
The second term is precisely the state-dependent correction to the coordinate correlation dynamics.

\section{Technical details for the examples}
\label{app:examples_details}

\subsection{CIR square-root diffusion benchmark}
\label{app:cir_benchmark_details}

We provide the analytic details for the Cox--Ingersoll--Ross benchmark used in Section~\ref{subsec:cir_benchmark}. Consider
\begin{equation}
    \d X_t
    =
    \kappa(\theta-X_t)\,\d t
    +
    \sqrt{2\gamma X_t}\,\d W_t,
    \qquad X_t>0,
    \label{eq:cir_sde_appendix}
\end{equation}
with \(\kappa,\theta,\gamma>0\). We introduce the dimensionless parameters
\begin{equation}
    \nu:=\frac{\kappa\theta}{\gamma},
    \qquad
    \beta:=\frac{\kappa}{\gamma}.
    \label{eq:cir_nu_beta_def}
\end{equation}
The stationary density is the Gamma distribution
\begin{equation}
    p_{\mathrm{ss}}(x)
    =
    \frac{\beta^\nu}{\Gamma(\nu)}
    x^{\nu-1}e^{-\beta x},
    \qquad x>0.
    \label{eq:cir_stationary_density_appendix}
\end{equation}
Equivalently,
\begin{equation}
    s(x)
    :=
    \partial_x\log p_{\mathrm{ss}}(x)
    =
    \frac{\nu-1}{x}-\beta
    =
    \frac{\kappa\theta-\gamma}{\gamma x}
    -
    \frac{\kappa}{\gamma}.
    \label{eq:cir_score_appendix}
\end{equation}

The corresponding Fokker--Planck equation is
\begin{equation}
    \partial_t p
    =
    -\partial_x\!\left[\kappa(\theta-x)p\right]
    +
    \partial_x^2\!\left[\gamma x p\right].
    \label{eq:cir_fokker_planck_appendix}
\end{equation}
At stationarity, imposing zero probability flux gives
\begin{equation}
    \kappa(\theta-x)p_{\mathrm{ss}}(x)
    -
    \partial_x\!\left[\gamma x p_{\mathrm{ss}}(x)\right]
    =
    0.
    \label{eq:cir_zero_flux_appendix}
\end{equation}
Dividing by \(p_{\mathrm{ss}}\) yields
\begin{equation}
    \kappa(\theta-x)
    =
    \gamma x\,\partial_x\log p_{\mathrm{ss}}(x)
    +
    \gamma.
    \label{eq:cir_score_drift_identity_appendix}
\end{equation}
Therefore the drift can be written as
\begin{equation}
    \kappa(\theta-x)
    =
    M(x)s(x)+\partial_x M(x),
    \qquad
    M(x)=\gamma x.
    \label{eq:cir_score_mobility_appendix}
\end{equation}
This gives the exact mobility decomposition
\begin{align}
    M(x)
    &=
    \Phi+\delta M(x),
    \\
    \Phi
    &=
    \langle M\rangle_{p_{\mathrm{ss}}}
    =
    \gamma\langle X\rangle_{p_{\mathrm{ss}}}
    =
    \gamma\theta,
    \\
    \delta M(x)
    &=
    \gamma(x-\theta).
    \label{eq:cir_exact_mobility_decomposition_appendix}
\end{align}

The exact transition density is known in terms of a modified Bessel function. Let
\begin{equation}
    z:=e^{-\kappa t},
    \qquad
    c_t:=\frac{\kappa}{\gamma(1-z)}=\frac{\beta}{1-z},
    \qquad
    u:=c_t z x_0,
    \qquad
    v:=c_t x.
    \label{eq:cir_transition_variables_appendix}
\end{equation}
Then
\begin{equation}
    p_t(x\mid x_0)
    =
    c_t
    \exp[-(u+v)]
    \left(\frac{v}{u}\right)^{(\nu-1)/2}
    I_{\nu-1}\!\left(2\sqrt{uv}\right),
    \qquad x>0,
    \label{eq:cir_transition_density_appendix}
\end{equation}
where \(I_q\) denotes the modified Bessel function of the first kind. In the original variables this reads
\begin{equation}
\begin{aligned}
    p_t(x\mid x_0)
    &=
    \frac{\kappa}{\gamma(1-e^{-\kappa t})}
    \exp\!\left[
        -\frac{\kappa(e^{-\kappa t}x_0+x)}
        {\gamma(1-e^{-\kappa t})}
    \right]
    \left(
        \frac{x}{e^{-\kappa t}x_0}
    \right)^{\frac{\kappa\theta/\gamma-1}{2}}
    \\
    &\hspace{2.5cm}\times
    I_{\kappa\theta/\gamma-1}\!\left(
        \frac{2\kappa\sqrt{e^{-\kappa t}x_0x}}
        {\gamma(1-e^{-\kappa t})}
    \right).
\end{aligned}
\label{eq:cir_transition_density_original_appendix}
\end{equation}

The conditional transition score used in the inverse identity is
\begin{equation}
    s_{t|0}(x\mid x_0)
    :=
    \partial_{x_0}\log p_t(x\mid x_0).
    \label{eq:cir_conditional_score_def_appendix}
\end{equation}
Starting from \eqref{eq:cir_transition_density_appendix},
\begin{equation}
\begin{aligned}
    \partial_{x_0}\log p_t(x\mid x_0)
    &=
    -\partial_{x_0}u
    -
    \frac{\nu-1}{2x_0}
    +
    \frac{I_{\nu-1}'(2\sqrt{uv})}{I_{\nu-1}(2\sqrt{uv})}
    \partial_{x_0}\!\left(2\sqrt{uv}\right).
\end{aligned}
\label{eq:cir_conditional_score_derivation_1}
\end{equation}
Since
\begin{equation}
    \partial_{x_0}u=c_t z,
    \qquad
    \partial_{x_0}\!\left(2\sqrt{uv}\right)
    =
    \frac{\sqrt{uv}}{x_0},
    \label{eq:cir_conditional_score_derivative_terms}
\end{equation}
and
\begin{equation}
    I_q'(r)
    =
    I_{q+1}(r)+\frac{q}{r}I_q(r),
    \label{eq:cir_bessel_derivative_identity}
\end{equation}
the singular terms cancel and we obtain
\begin{equation}
    s_{t|0}(x\mid x_0)
    =
    c_t z
    \left[
        -1
        +
        \sqrt{\frac{x}{z x_0}}\,
        \frac{
            I_{\nu}\!\left(2c_t\sqrt{z x_0 x}\right)
        }{
            I_{\nu-1}\!\left(2c_t\sqrt{z x_0 x}\right)
        }
    \right].
    \label{eq:cir_conditional_score_appendix}
\end{equation}

We now compute the lagged correlations. For \(\phi_\alpha(x)=x^\alpha\), the conditional moment is
\begin{equation}
    m_\alpha(t,x_0)
    :=
    \mathbb{E}[X_t^\alpha\mid X_0=x_0]
    =
    c_t^{-\alpha}
    \frac{\Gamma(\alpha+\nu)}{\Gamma(\nu)}
    {}_1F_1\!\left(
        -\alpha;\nu;-c_t z x_0
    \right),
    \label{eq:cir_conditional_moment_appendix}
\end{equation}
where \({}_1F_1\) is Kummer's confluent hypergeometric function. Averaging against the stationary density gives
\begin{equation}
\begin{aligned}
    C_{\alpha,1}(t)
    &:=
    \mathbb{E}_{p_{\mathrm{ss}}}
    \left[
        X_t^\alpha X_0
    \right]
    \\
    &=
    \int_0^\infty
    m_\alpha(t,x_0)\,
    x_0\,
    p_{\mathrm{ss}}(x_0)\,\d x_0
    \\
    &=
    \beta^{-(\alpha+1)}
    \frac{\Gamma(\alpha+\nu)}{\Gamma(\nu)}
    \left(
        \nu+\alpha z
    \right).
\end{aligned}
\label{eq:cir_power_coordinate_correlation_appendix}
\end{equation}
Therefore
\begin{equation}
    \dot C_{\alpha,1}(t)
    =
    -\kappa\alpha z\,
    \beta^{-(\alpha+1)}
    \frac{\Gamma(\alpha+\nu)}{\Gamma(\nu)}.
    \label{eq:cir_correlation_derivative_beta_appendix}
\end{equation}
Using \(\nu=\kappa\theta/\gamma\), \(\beta=\kappa/\gamma\), and \(z=e^{-\kappa t}\), this is
\begin{equation}
    \dot C_{\alpha,1}(t)
    =
    -\alpha\theta\gamma
    \left(\frac{\kappa}{\gamma}\right)^{1-\alpha}
    \frac{
        \Gamma(\alpha+\kappa\theta/\gamma)
    }{
        \Gamma(\kappa\theta/\gamma+1)
    }
    e^{-\kappa t}.
    \label{eq:cir_correlation_derivative_appendix}
\end{equation}

We next evaluate the two sides of the inverse identity for the affine family
\begin{equation}
    \delta M_a(x)=a(x-\theta).
    \label{eq:cir_affine_family_appendix}
\end{equation}
First,
\begin{equation}
    \partial_{x_0}m_\alpha(t,x_0)
    =
    \alpha z\,c_t^{1-\alpha}
    \frac{\Gamma(\alpha+\nu)}{\Gamma(\nu+1)}
    {}_1F_1\!\left(
        1-\alpha;\nu+1;-c_t z x_0
    \right).
    \label{eq:cir_derivative_conditional_moment_appendix}
\end{equation}
Since
\begin{equation}
    \partial_{x_0}m_\alpha(t,x_0)
    =
    \mathbb{E}\!\left[
        X_t^\alpha
        s_{t|0}(X_t\mid x_0)
        \,\middle|\,
        X_0=x_0
    \right],
    \label{eq:cir_conditional_score_moment_identity_appendix}
\end{equation}
the operator acting on \(\delta M_a\) is
\begin{equation}
\begin{aligned}
    \mathcal{A}_{t,\alpha,1}[\delta M_a]
    &:=
    -
    \left\langle
        X_t^\alpha
        s_{t|0}(X_t\mid X_0)
        \delta M_a(X_0)
    \right\rangle_{\mathrm{obs}}
    \\
    &=
    a
    \left\langle
        (X_0-\theta)
        \partial_{x_0}m_\alpha(t,X_0)
    \right\rangle_{p_{\mathrm{ss}}}.
\end{aligned}
\label{eq:cir_operator_affine_appendix}
\end{equation}
Using the Gamma average of \({}_1F_1\), this gives
\begin{equation}
    \mathcal{A}_{t,\alpha,1}[\delta M_a]
    =
    -a\,K_\alpha(t),
    \label{eq:cir_operator_affine_reduction_appendix}
\end{equation}
where
\begin{equation}
    K_\alpha(t)
    :=
    \alpha\theta z
    \beta^{1-\alpha}
    \frac{\Gamma(\alpha+\nu)}{\Gamma(\nu+1)}
    \left[
        1
        -
        {}_2F_1\!\left(
            1-\alpha,1;\nu+1;z
        \right)
    \right].
    \label{eq:cir_K_alpha_def_appendix}
\end{equation}
Here \({}_2F_1\) denotes the Gauss hypergeometric function.

The observed residual appearing in the split identity is
\begin{equation}
    \mathcal{E}_{\alpha,1,\mathrm{obs}}(t)
    =
    \dot C_{\alpha,1,\mathrm{obs}}(t)
    -
    \Phi
    \left\langle
        X_t^\alpha s(X_0)
    \right\rangle_{\mathrm{obs}}
    \label{eq:cir_residual_def_appendix}
\end{equation}
By integration by parts with respect to the stationary density,
\begin{equation}
    \left\langle
        X_t^\alpha s(X_0)
    \right\rangle_{\mathrm{obs}}
    =
    -
    \left\langle
        \partial_{x_0}m_\alpha(t,X_0)
    \right\rangle_{p_{\mathrm{ss}}}.
    \label{eq:cir_score_corr_integration_by_parts_appendix}
\end{equation}
Using \eqref{eq:cir_derivative_conditional_moment_appendix}, one obtains
\begin{equation}
    \left\langle
        X_t^\alpha s(X_0)
    \right\rangle_{\mathrm{obs}}
    =
    -
    \alpha z
    \beta^{1-\alpha}
    \frac{\Gamma(\alpha+\nu)}{\Gamma(\nu+1)}
    {}_2F_1\!\left(
        1-\alpha,1;\nu+1;z
    \right).
    \label{eq:cir_score_correlation_appendix}
\end{equation}
Since \(\Phi=\theta\gamma\), combining
\eqref{eq:cir_correlation_derivative_appendix}
and
\eqref{eq:cir_score_correlation_appendix}
gives
\begin{equation}
    \mathcal{E}_{\alpha,1,\mathrm{obs}}(t)
    =
    -\gamma K_\alpha(t).
    \label{eq:cir_residual_reduction_appendix}
\end{equation}
Therefore the inverse equation
\begin{equation}
    \mathcal{A}_{t,\alpha,1}[\delta M_a]
    =
    \mathcal{E}_{\alpha,1,\mathrm{obs}}(t)
    \label{eq:cir_inverse_equation_appendix}
\end{equation}
reduces to
\begin{equation}
    aK_\alpha(t)=\gamma K_\alpha(t).
    \label{eq:cir_inverse_scalar_equation_appendix}
\end{equation}

For the coordinate observable, \(\alpha=1\), one has
\begin{equation}
    {}_2F_1(0,1;\nu+1;z)=1,
    \qquad
    K_1(t)=0.
    \label{eq:cir_coordinate_degeneracy_appendix}
\end{equation}
Thus the coordinate channel cannot identify the affine correction. Equivalently,
\begin{equation}
    \mathbb{E}[X_t\mid X_0=x_0]
    =
    \theta+z(x_0-\theta),
    \label{eq:cir_affine_conditional_mean_appendix}
\end{equation}
so that
\begin{equation}
\begin{aligned}
    \left\langle
        X_t s_{t|0}(X_t\mid X_0)\delta M_a(X_0)
    \right\rangle_{\mathrm{obs}}
    &=
    \left\langle
        \partial_{x_0}\mathbb{E}[X_t\mid X_0]\,
        a(X_0-\theta)
    \right\rangle_{p_{\mathrm{ss}}}
    \\
    &=
    az
    \left\langle
        X_0-\theta
    \right\rangle_{p_{\mathrm{ss}}}
    =
    0.
\end{aligned}
\label{eq:cir_coordinate_nullspace_appendix}
\end{equation}

For any nonlinear power observable \(\alpha\neq 1\), the factor \(K_\alpha(t)\) is not identically zero as a function of the lag. Indeed, its small-\(z\) expansion is
\begin{equation}
    1
    -
    {}_2F_1(1-\alpha,1;\nu+1;z)
    =
    \frac{\alpha-1}{\nu+1}z
    +
    \mathcal{O}(z^2),
    \qquad z\to 0.
    \label{eq:cir_K_nonzero_expansion_appendix}
\end{equation}
Thus the affine coefficient is uniquely identified:
\begin{equation}
    a=\gamma.
    \label{eq:cir_a_gamma_appendix}
\end{equation}
Consequently,
\begin{equation}
    \delta M(x)=\gamma(x-\theta),
    \qquad
    M(x)=\theta\gamma+\gamma(x-\theta)=\gamma x.
    \label{eq:cir_final_mobility_appendix}
\end{equation}
This proves exact recovery of the CIR mobility from the correlation constraints once the observable library contains a nonlinear observable that removes the coordinate-channel nullspace.


\subsection{Technical details for the two-dimensional affine multiplicative-noise benchmark}
\label{app:affine_multiplicative_2d_details}

This appendix gives the experiment-specific implementation details behind the
two-dimensional benchmark reported in
Section~\ref{subsec:affine_multiplicative_2d_results}. The reference SDE,
its coefficients, the diagnostic construction of
\(\bm{M}_{\mathrm{ref}}\), and the two forward-validation Langevin equations are
stated in the main text. Here we collect the remaining details needed to build
the empirical residual targets and to train the neural mobility correction.

The stationary and conditional scores are estimated by the denoising
score-matching constructions of Appendix~\ref{app:score_estimation_dsm}. In
the present experiment we use perturbation level \(\sigma=0.05\), a
stationary-score network with hidden widths \((128,64)\), and a lag-conditioned
joint-score network with hidden widths \((256,128)\).

\subsubsection{Empirical correlation residuals}

Let \(\mathcal{T}=\{\tau_1,\dots,\tau_K\}\subset(0,\infty)\) denote the discrete positive lag grid used throughout the conditional-score fit and the mobility fit. Correlations are evaluated on \(\{0\}\cup\mathcal{T}\), with \(\tau=0\) entering only through the baseline estimate for \(\bm{\Phi}\).

For observables \(\phi_m,\phi_n\), the empirical lagged cross-correlation is computed from all admissible lagged pairs:
\begin{equation}
    \widehat{\bm{C}}_{m,n}(\tau)
    =
    \frac{1}{N_\tau}
    \sum_{r=1}^{N_\tau}
    \phi_m(\bm{x}^{(r)}_\tau)
    \phi_n(\bm{x}^{(r)}_0)^T.
    \label{eq:app_affine_2d_empirical_correlation}
\end{equation}
For each matrix entry of \(\widehat{\bm{C}}_{m,n}(\tau)\), we fit an entrywise cubic interpolant in \(\tau\) on the grid \(\{0\}\cup\mathcal{T}\). The derivative \(\widehat{\dot{\bm{C}}}_{m,n,\mathrm{obs}}(\tau)\) is then the analytic derivative of this interpolant evaluated at the positive lags \(\tau\in\mathcal{T}\), and the baseline mobility is estimated from the corresponding one-sided right derivative at the origin:
\begin{equation}
    \bm{\Phi}
    =
    -
    \widehat{\dot{\bm{C}}}_{1,1,\mathrm{obs}}(0^+),
    \qquad
    \phi_1(\bm{x})=\bm{x}.
    \label{eq:app_affine_2d_phi_estimator}
\end{equation}
For each observable pair and component \(j\), we compute the empirical
approximation of \(B_{\bm{b}_{\bm{\Phi},n,j}}\), where
\(\bm{b}_{\bm{\Phi},n,j}=-\bm{\Phi}\nabla\phi_{n,j}\),
\begin{equation}
    \widehat{B}_{\bm{b}_{\bm{\Phi},n,j}}(\bm{x})
    =
    \nabla\cdot
    \left(
        \bm{\Phi}\nabla\phi_{n,j}(\bm{x})
    \right)
    +
    \nabla\phi_{n,j}(\bm{x})^T\bm{\Phi}^T
    \widehat{\bm{s}}_\psi(\bm{x}),
    \label{eq:app_affine_2d_gamma_hat}
\end{equation}
and the empirical baseline response term
\begin{equation}
    \left[\widehat{\bm{G}}_{m,n}(\tau)\right]_{:,j}
    =
    \frac{1}{N_\tau}
    \sum_{r=1}^{N_\tau}
    \phi_m(\bm{x}^{(r)}_\tau)
    \widehat{B}_{\bm{b}_{\bm{\Phi},n,j}}(\bm{x}^{(r)}_0).
    \label{eq:app_affine_2d_G_hat}
\end{equation}
The data-derived residual response is then
\begin{equation}
    \widehat{\bm{\mathcal E}}_{m,n,\mathrm{obs}}(\tau)
    =
    \widehat{\dot{\bm{C}}}_{m,n,\mathrm{obs}}(\tau)
    -
    \widehat{\bm{G}}_{m,n}(\tau)
    \label{eq:app_affine_2d_E_hat}
\end{equation}
This is the GFDT response channel that the learned mobility correction must
reproduce. Equivalently, it is the part of the correlation derivative not
explained by the constant mean-mobility impulse response.

The observable channels used in the experiment are
\begin{equation}
    \phi_n\in\{x,y\},
    \qquad
    \phi_m\in
    \{x,y,x^2,xy,y^2,x^3,x^2y,xy^2,y^3\}.
    \label{eq:app_affine_2d_observable_library}
\end{equation}
This gives \(18\) ordered channels. The coordinate observables determine the baseline response sector, while the quadratic and cubic observables probe state-dependent departures from the constant-mobility closure.

\subsubsection{Mobility parameterization and training objective}

The learned mobility is written as
\begin{equation}
    \bm{M}_\theta(\bm{x})
    =
    \bm{D}_\theta(\bm{x})
    +
    \bm{R}_\theta(\bm{x}),
    \qquad
    \delta\bm{M}_\theta(\bm{x})
    =
    \bm{M}_\theta(\bm{x})-\bm{\Phi}.
    \label{eq:app_affine_2d_M_theta_split}
\end{equation}
In two dimensions, the efficient parameterization uses three scalar outputs for the symmetric diffusion tensor and one scalar output for the antisymmetric circulation. Specifically,
\begin{equation}
    \bm{D}_\theta(\bm{x})
    =
    \bm{L}_\theta(\bm{x})
    \bm{L}_\theta(\bm{x})^T
    +
    \varepsilon\bm{I},
    \qquad
    \bm{L}_\theta(\bm{x})
    =
    \begin{pmatrix}
        \ell_{11,\theta}(\bm{x}) & 0\\
        \ell_{21,\theta}(\bm{x}) & \ell_{22,\theta}(\bm{x})
    \end{pmatrix},
    \qquad
    \varepsilon>0,
    \label{eq:app_affine_2d_D_parameterization}
\end{equation}
and
\begin{equation}
    \bm{R}_\theta(\bm{x})
    =
    r_\theta(\bm{x})
    \begin{pmatrix}
        0 & -1\\
        1 & 0
    \end{pmatrix}.
    \label{eq:app_affine_2d_R_parameterization}
\end{equation}
Thus the network outputs only the four independent scalar fields
\[
    \ell_{11,\theta},\quad
    \ell_{21,\theta},\quad
    \ell_{22,\theta},\quad
    r_\theta,
\]
rather than a redundant unconstrained \(2\times2\) matrix. This enforces positive semidefiniteness of the symmetric part by construction and imposes antisymmetry of the circulation exactly.

For a trial mobility correction, the predicted residual response is evaluated in
the equivalent conditional-score form as
\begin{align}
    \widehat{\bm{\mathcal A}}_{m,n}^{\theta}(\tau)
    &=
    -
    \frac{1}{B}
    \sum_{r=1}^{B}
    \phi_m(\bm{x}^{(r)}_\tau)\,
    \widehat{\bm{s}}_{\tau|0}(\bm{x}^{(r)}_\tau\mid\bm{x}^{(r)}_0)^T
    \\
    &\quad\cdot
    \delta\bm{M}_\theta(\bm{x}^{(r)}_0)\,
    \nabla\phi_n(\bm{x}^{(r)}_0)^T.
    \label{eq:app_affine_2d_B_theta}
\end{align}
The mobility loss is
\begin{equation}
\begin{aligned}
    \mathcal{L}_{\mathrm{mob}}(\theta)
    &=
    \sum_{(m,n)\in\mathcal{I}_\phi}
    \sum_{\tau\in\mathcal{T}}
    w_\tau
    \left\|
        \frac{
            \widehat{\bm{\mathcal E}}_{m,n,\mathrm{obs}}(\tau)
            -
            \widehat{\bm{\mathcal A}}_{m,n}^{\theta}(\tau)
        }{
            S_{m,n}
        }
    \right\|_F^2 \\
    &\quad
    +
    \lambda_{\mathrm{mean}}
    \left\|
        \frac{1}{B_{\mathrm{anc}}}
        \sum_{r=1}^{B_{\mathrm{anc}}}
        \delta\bm{M}_\theta(\bm{x}^{(r)}_{\mathrm{anc}})
    \right\|_F^2
    +
    \lambda_{\mathrm{reg}}\mathcal{R}(\theta).
\end{aligned}
    \label{eq:app_affine_2d_mobility_loss}
\end{equation}
Here \(\mathcal{I}_\phi\) is the set of the \(18\) fitted observable channels, \(S_{m,n}\) is the empirical Frobenius-norm scale of the target residual curve for channel \((m,n)\), and \(w_\tau\) is a prescribed positive lag weight that mildly emphasizes the shortest resolved lags. The mean penalty enforces the decomposition
\[
    \left\langle
        \delta\bm{M}_\theta
    \right\rangle_{\mathrm{obs}}
    \approx
    \bm{0}
\]
without subtracting a mini-batch mean from \(\bm{D}_\theta\), which would generally destroy the positive-semidefinite parameterization.

The mobility network has two hidden layers of width \((128,128)\). It is initialized at the constant baseline and trained for \(250\) Adam epochs with learning rate \(3\times10^{-4}\), weight decay \(10^{-6}\), \(32768\) sampled lagged pairs per lag, mini-batches of four lag values, and \(32768\) anchor states for the mean-correction penalty.

\begin{figure*}[p]
    \centering
    \includegraphics[width=\textwidth,height=0.88\textheight,keepaspectratio]{fit_A.png}
    \caption{Training residual channels for the two-dimensional affine multiplicative-noise benchmark. The black curves are the data-derived residual targets, the orange dashed curves are ex-post evaluations of the same empirical operator using the reference mobility correction, and the blue dash-dotted curves are the learned network predictions.}
    \label{fig:affine_A_fit}
\end{figure*}

Figure~\ref{fig:affine_A_fit} shows the objects directly optimized during mobility training. Across the \(18\) fitted channels, the learned operator reproduces the scale and lag dependence of the data-derived residuals, including the rapid decay of the quadratic and cubic channels and the sign-changing behavior in the \(B_{y,x}\) channel. The largest discrepancies occur in several coordinate and mixed channels at the shortest lags, where the target is most sensitive to the derivative of the empirical correlation curve and to conditional-score error. The mean RMSE between the data-derived targets and the learned operator over the fitted channels is \(1.64\times10^{-3}\), the mean normalized RMSE is \(4.05\times10^{-1}\), and the mean RMSE against the ex-post reference-operator curves is \(1.07\times10^{-3}\).

\bibliographystyle{unsrtnat}
\bibliography{references}
\end{document}
